\chapter{多模态理解}
\label{chap:vision-multimodal-understanding}

给一张街景照片，要求找出“右边拿杯子的人”，结果应落到一个可检查的区域。给一页票据，询问“总额是多少”，模型需要读出数字及其字段归属。给一段视频，追问“第二次敲门发生在讲话之前还是之后”，还需对齐事件、声音与时间。三个请求都可以用一句话回答，但支持这句话的证据并不相同：第一个需要区域，第二个需要字符串和字段关系，第三个需要带时间的事件记录。

本章要回答的问题是：怎样把文字、图像、视频、声音与页面结构联系起来，找到支持结果的证据，并给出可以复核的描述和答案。\ref{chap:vision-multimodal-alignment}章已经讲过跨模态对齐与融合的基本原理，包括典型相关分析、动态时间规整、注意力对齐、对比对齐以及早期、中间、晚期三种融合位置；\ref{chap:vision-image-features}章讲过边缘、局部特征、阈值分割和模板匹配等经典图像特征；\ref{chap:vision-multimodal-models}章讲过多模态模型怎样匹配表示、连接模态和解码语言；\ref{chap:vision-image-recognition}章和\ref{chap:vision-video-recognition}章讲过图像与视频中的检测、跟踪和检索。本章使用这些原理和接口，按请求所需的成果展开，用到时只回顾结论并指回原处。

路线由证据从粗到细的获取顺序决定。第一节先规定任务的输入、输出与评价指标，并说明多模态理解在对齐粒度、监督质量、语言先验、观测分辨率和证据组合上的困难。面对大量图片或视频，先要缩小候选范围，这是跨模态匹配与检索；候选找到后，要把查询落到区域、时间或控件上，这是跨模态定位；票据、表格和图表的内容写在像素里，还需读取文字并恢复结构；有了区域、文字和事件，才能选择内容形成描述，或者组合证据回答问题。最后一节讨论怎样分层核验这些结果。贯穿全章的是来源记录：每个判断依据哪幅图、哪一页、哪个区域或哪段时间。输入额外的OCR结果、字幕或真实数据表，会改变模型仍需完成的工作，因而也改变评价结论。

\section{多模态理解的问题与挑战}
\label{sec:vision-understanding-challenges}

开篇的三个请求有一个共同结构：系统收到一份或多份观测，例如图像、视频帧、音频或页面截图，同时收到一个用另一种模态表达的请求，例如一句话、一段声音或一幅示例图；它要输出一个能回到观测中检查的结果。按结果的形式，多模态理解的输出可以归为四类。第一类是排序列表，即从候选库中挑出相关项并给出名次；第二类是位置，即框、掩码、点击点或时间区间；第三类是字符串与结构，即读出的文字、表格的行列、表单的键值、图表的数据和图解的连接；第四类是自然语言，即描述和答案。不同输出能被检查的方式不同，评价指标也随之不同。后文各节比较模型时都要固定一种指标，因此这里先把指标定义清楚。

排序列表用召回类指标评价。设查询集合为$Q$，对查询$q$，相关候选的集合为$R_q$，模型排在前$K$名的候选集合为$S_q^K$。图文检索中最常用的是前$K$召回率（Recall@$K$，记作R@$K$）：
\begin{equation}
  \mathrm{R@}K=\frac{1}{|Q|}\sum_{q\in Q}
    \bm 1\bigl[S_q^K\cap R_q\neq\varnothing\bigr]\text{，}
  \label{eq:vision-understanding-recall-at-k}
\end{equation}
其中$\bm 1[\cdot]$为指示函数，条件成立取1，否则取0。它统计“前$K$名里至少出现一个正例”的查询比例。一个查询有多个正例时，也可以改用覆盖比例$|S_q^K\cap R_q|/|R_q|$，两种定义在多对多数据上给出不同数值，报告时必须写明。另一个常用量是\emph{中位名次}（median rank, MedR）：每个查询第一个正例的名次取中位数，越小越好。以三个查询为教学例，若第一个正例分别排在第1、4、12名，则R@1$=1/3$，R@5$=2/3$，R@10$=2/3$，MedR$=4$。R@$K$只关心正例是否进入前$K$名，名次在前$K$之内如何排列并不影响它，这正是后文把检索拆成召回与重排的依据。

位置输出用重叠度评价。框$B$与标注框$B^*$的\term{交并比}（intersection over union, IoU）定义为
\begin{equation}
  \operatorname{IoU}(B,B^*)=\frac{|B\cap B^*|}{|B\cup B^*|}\text{，}
  \label{eq:vision-understanding-iou}
\end{equation}
$|\cdot|$表示面积。指代表达定位通常报告IoU不小于0.5的样本比例，即Acc@0.5。时间区间用同一公式，只是把面积换成长度，称为时间IoU。例如框$[0,4]\times[0,2]$与标注框$[2,6]\times[0,2]$的交集面积为4，并集面积为12，IoU$=1/3$，按0.5的阈值记为错误；区间$[1,3)$与$[2,4)$的时间IoU同样是$1/3$。界面定位输出的是点击点，常用的口径是点是否落在目标控件框内；声源定位输出热图，要先按阈值取出区域，再与标注的声源区域比较。阈值本身是协议的一部分：同一组输出，在0.5与0.7两个阈值下可以给出差别很大的准确率。

字符串用编辑距离评价。两个字符串之间的\term{编辑距离}（edit distance），又称Levenshtein距离，是把一个串变成另一个串所需的最少替换、删除和插入次数。设参考串$\bm a=a_1\cdots a_n$与识别串$\bm b=b_1\cdots b_m$，记$d_{i,j}$为$a_1\cdots a_i$与$b_1\cdots b_j$之间的编辑距离。最后一步只有三种可能：删去$a_i$、插入$b_j$，或者让$a_i$与$b_j$对应（相同则不计代价，不同则计一次替换），于是
\begin{equation}
  d_{i,j}=\min\bigl\{d_{i-1,j}+1,\ d_{i,j-1}+1,\
    d_{i-1,j-1}+\bm 1[a_i\neq b_j]\bigr\}\text{，}
  \qquad d_{i,0}=i,\ d_{0,j}=j\text{。}
  \label{eq:vision-understanding-edit-distance}
\end{equation}
边界条件说明空串与长度为$i$的串之间需要$i$次插入或删除。这个递推是动态规划：$d_{i,j}$只依赖左、上、左上三个已经算好的格，按行填表即可，代价为$O(nm)$。动态规划把问题分解为重叠子问题、按依赖顺序求解的原理，可参考\BookChapterRef{chap:graph-theory}{卷一“图的表示与计算”章}中的最短路递推；编辑距离的递推相当于在$(n+1)\times(m+1)$的网格上，从左上角走到右下角的最短路径，横、竖每步代价为1，斜步代价为0或1。算法\ref{alg:vision-understanding-edit-distance}写出了计算步骤。

\begin{algorithm}[字符串编辑距离的动态规划计算]
  \label{alg:vision-understanding-edit-distance}
  \begin{algorithmic}[1]
    \Require 参考串$a_1\cdots a_n$，识别串$b_1\cdots b_m$
    \Ensure 编辑距离$d_{n,m}$
    \For{$i=0,1,\ldots,n$}
      \State $d_{i,0}\gets i$
    \EndFor
    \For{$j=0,1,\ldots,m$}
      \State $d_{0,j}\gets j$
    \EndFor
    \For{$i=1,\ldots,n$}
      \For{$j=1,\ldots,m$}
        \State 按式\eqref{eq:vision-understanding-edit-distance}取三种来源的最小值作为$d_{i,j}$
      \EndFor
    \EndFor
    \State \Return $d_{n,m}$；需要具体的编辑操作时，从$(n,m)$沿取得最小值的来源回溯到$(0,0)$
  \end{algorithmic}
\end{algorithm}

以教学票据编号为例，参考串为“A017”，识别器把数字0读成字母O并漏掉了1，输出“AO7”。表\ref{tab:vision-understanding-edit-example}按行填出$d_{i,j}$。右下角为2，对应一次替换（0读成O）和一次删除（漏掉1）。

\begin{table}[htbp]
  \centering
  \small
  \begin{tabular}{@{}lcccc@{}}
    \toprule
     & 空 & A & O & 7\\
    \midrule
    空 & 0 & 1 & 2 & 3\\
    A & 1 & 0 & 1 & 2\\
    0 & 2 & 1 & 1 & 2\\
    1 & 3 & 2 & 2 & 2\\
    7 & 4 & 3 & 3 & 2\\
    \bottomrule
  \end{tabular}
  \caption{参考串“A017”（行）与识别串“AO7”（列）的编辑距离表，为教学构造。每格取上方加1、左方加1、左上方加替换代价三者的最小值；例如第“0”行“O”列取左上方的0加1，得到1。右下角的2是两串的编辑距离。}
  \label{tab:vision-understanding-edit-example}
\end{table}

用编辑距离可以定义\term{字符错误率}（character error rate, CER）。若参考串含$N>0$个字符，最优编辑中的替换、删除和插入次数分别为$S,D,I$，则
\begin{equation}
  \mathrm{CER}=\frac{S+D+I}{N}\text{，}
  \label{eq:vision-cer}
\end{equation}
分子就是编辑距离。上例中CER$=2/4=0.5$。分母使用参考字符数，因此插入很多字符时CER可以超过1；空格、标点和大小写是否计入也要事先规定。以词为单位做同样的计算，得到词错误率。

答案中的字符串还有一种更宽容的指标。场景文字问答与文档问答常用\term{平均归一化编辑相似度}（average normalized Levenshtein similarity, ANLS）\scite{vision-stvqa2019}。设第$i$个问题有若干参考答案$a_{i,j}$，模型答案为$o_i$，归一化编辑距离$\operatorname{NL}$等于编辑距离除以两串中较长者的长度，取值在0到1之间。定义
\begin{equation}
  \mathrm{ANLS}=\frac1N\sum_{i=1}^{N}\max_j s(a_{i,j},o_i)\text{，}\qquad
  s(a,o)=\begin{cases}
    1-\operatorname{NL}(a,o), & \operatorname{NL}(a,o)<\tau,\\
    0, & \operatorname{NL}(a,o)\geq\tau,
  \end{cases}
  \label{eq:vision-understanding-anls}
\end{equation}
原文取$\tau=0.5$。阈值的用意是：错了一半以上的字符串更可能选错了文字，而不是读错了个别字符，此时不给部分分。例如参考答案“128.50”，回答“128.5”的编辑距离为1，$\operatorname{NL}=1/6$，得分约0.83；回答“100.00”需要三次替换，$\operatorname{NL}=3/6=0.5$，得分为0。ANLS宽容OCR的小错，但对金额这样的字段，少一位小数可能已经改变含义，因此仍要配合整串正确率。

表格结构需要比较两张表的行列和跨格关系，而不只是比较字符串。\term{树编辑距离相似度}（tree-edit-distance-based similarity, TEDS）把表格写成HTML树，计算
\begin{equation}
  \operatorname{TEDS}(T_a,T_b)=1-\frac{\operatorname{EditDist}(T_a,T_b)}{\max(|T_a|,|T_b|)}\text{，}
  \label{eq:vision-understanding-teds}
\end{equation}
其中$\operatorname{EditDist}$为把一棵树变成另一棵所需的最少节点插入、删除和替换代价，$|T|$为节点数\scite{vision-pubtabnet2020}。它把字符串编辑距离推广到树：节点是表格、行和单元格，单元格内容不同也计入替换代价。PubTables-1M另用GriTS系列指标，把表格看成单元格矩阵，在两个矩阵最相似的行列子结构上计算F分数形式的相似度，并分别考察拓扑、内容和位置\scite{vision-pubtables2022}。页面版面分析则把标题、正文、表格等区域当作检测对象，用\term{平均精度}（average precision, AP）评价。

视觉问答的答案通常很短。VQA数据集为每个问题收集10个人工答案，单个回答的准确率定义为$\min(\text{给出该答案的人数}/3,\,1)$，并在10个答案中任取9个形成的10个子集上取平均\scite{vision-vqa2015}。若10人中只有2人回答“蓝色”，模型也回答“蓝色”：在包含这2人的8个子集上得分$2/3$，在只含其中1人的2个子集上得分$1/3$，平均为$(8\times2/3+2\times1/3)/10=0.6$。这个定义承认答案的多样性，三人以上同意即得满分；但它只比较答案字符串，不检查答案依据的是图中哪一部分。图表问答中的数值答案常用放宽准确率：与正确值相差不超过5\%即算正确\scite{vision-chartqa2022}。

描述的评价要比较一句生成文字与若干参考描述。BLEU计算候选句$n$元组在参考句中出现的比例，即$n$元精确率$p_n$，同一个$n$元组的计数不超过它在任一参考句中出现的最多次数，以免重复一个高频词就得到高分；再用简短惩罚$\mathrm{BP}$抑制过短的句子\scite{papineni-etal-2002-bleu}：
\begin{equation}
  \mathrm{BLEU}=\mathrm{BP}\cdot\exp\Bigl(\sum_{n=1}^{N}w_n\log p_n\Bigr)\text{，}\qquad
  \mathrm{BP}=\begin{cases}1, & c>r,\\ e^{1-r/c}, & c\leq r,\end{cases}
  \label{eq:vision-understanding-bleu}
\end{equation}
其中$c$为候选长度，$r$为参考长度，常用$N=4$与$w_n=1/4$。CIDEr针对图像描述设计：它先把每个$n$元组按TF-IDF加权，即该$n$元组在句中出现越多权重越高，在整个数据集的描述中越普遍权重越低；再计算候选句与每句参考描述的加权向量余弦相似度并取平均，最后对$n=1,\ldots,4$等权求和\scite{vision-cider2015}。设$\bm g_n(c)$为候选句$c$在长度为$n$的全部$n$元组上的TF-IDF向量，$s_1,\ldots,s_m$为参考描述，则
\begin{equation}
 \begin{gathered}
  \mathrm{CIDEr}_n(c)=\frac1m\sum_{j=1}^{m}
    \frac{\bm g_n(c)^\top\bm g_n(s_j)}{\lVert\bm g_n(c)\rVert_2\,\lVert\bm g_n(s_j)\rVert_2}\text{，}\\
  \mathrm{CIDEr}(c)=\sum_{n=1}^{4}\frac14\,\mathrm{CIDEr}_n(c)\text{。}
 \end{gathered}
  \label{eq:vision-understanding-cider}
\end{equation}
IDF权重压低了“一个”“在”这类在各图描述中都会出现的词组，突出与当前图像相关的内容词。常用的CIDEr-D版本另加长度惩罚和计数截断，并把分数乘以10。两种指标衡量的都是与参考描述的措辞重合，并不直接检查描述中的对象是否真的出现在图中；参考描述没有提到但图中确实存在的内容，同样得不到分数。

表\ref{tab:vision-understanding-metrics}汇总这些指标及其检查范围。每个指标只检查一种成果，任务需要同时满足几种条件时，要分别报告，或者按\ref{sec:vision-understanding-grounded-evaluation}节的方式取交集。

\begin{table}[htbp]
  \centering
  \small
  \begin{tabularx}{\linewidth}{@{}p{14mm}p{20mm}XX@{}}
    \toprule
    输出 & 指标 & 检查内容 & 不能检查的内容\\
    \midrule
    排序列表 & R@$K$、MedR & 正例是否进入前$K$名、第一个正例的名次 & 正例在前$K$名内的顺序；关系是否由同一对象满足\\
    位置 & IoU、Acc@0.5、点命中 & 区域或区间与标注的重叠 & 选中的是否为查询所指对象之外的同类对象\\
    字符串 & CER、整串正确率 & 字符级差异 & 字符串属于哪个字段、单位是否正确\\
    表格与页面 & TEDS、GriTS、AP & 行列、跨格、区域类别 & 内容的数值运算是否正确\\
    答案 & VQA准确率、ANLS、放宽准确率 & 答案字符串或数值 & 答案依据的区域和时刻\\
    描述 & BLEU、CIDEr & 与参考描述的$n$元组重合 & 描述中的对象和关系是否有图像依据\\
    \bottomrule
  \end{tabularx}
  \caption{多模态理解的主要评价指标。中间一列是指标直接计算的内容，最后一列是同一指标无法发现的错误，需要另设检查。}
  \label{tab:vision-understanding-metrics}
\end{table}

指标定义了“做对”的标准，下面各小节讨论为什么做对很难。困难来自多模态数据本身：不同模态的对应关系有不同粒度，训练配对往往不精确，语言一侧的先验很强，视觉一侧的证据又可能小而分散，还常常需要多步组合。每个困难都有经典方法的处理方式，也有学习方法带来的改变。

\subsection{多粒度对应带来的对齐困难}
\label{sec:vision-understanding-challenge-granularity}

“穿红衣服的人正在骑车”与一幅街景之间的对应，可以在几个粒度上成立。整句与整幅图可以对应，说明这幅图大体相关；短语“红衣服的人”与图中某个人物区域对应，说明属性绑定到了哪个对象；视频中“搬起”这个词与若干帧对应，说明动作发生在哪段时间；敲门声与手碰到门的那几帧对应，说明声音和画面同步。这些粒度不能相互替代：整图与整句的匹配分数再高，也推不出红色属于骑车的人；短语与区域对齐了，也推不出动作发生在哪一秒。困难在于，系统必须在任务所需的粒度上保留对应关系，而粗粒度的表示往往已经丢掉了细粒度对应所需的位置和时间信息。

经典方法把对齐写成显式的优化问题，见\ref{chap:vision-multimodal-alignment}章。典型相关分析（\ref{sec:vision-alignment-cca}节）为两组成对的特征各学一个线性投影，使投影后的相关系数最大，从而得到一个公共空间；它处理的是整体粒度，每个样本只有一个向量，无法表达区域与短语的对应。动态时间规整（\ref{sec:vision-alignment-dtw}节）在单调和端点约束下，用动态规划找出两条序列之间累计代价最小的对应路径，适合字幕与帧、音频与视频这类保持先后顺序的序列；但它要求事先给出逐帧之间的距离，距离本身就要求两种特征可比。这些方法所用的视觉特征来自\ref{chap:vision-image-features}章的人工设计，例如HOG、SIFT和视觉词袋，它们描述外观，却不知道哪个外观对应哪个词。

学习方法改变的是两点。其一，特征和对齐一起学习：对比对齐（\ref{sec:vision-alignment-contrastive}节）在大规模配对数据上同时训练两个编码器，让配对的表示彼此靠近、不配对的彼此远离，可比性由数据学出，不再依赖人工设计的特征。其二，细粒度对应可以通过注意力表达（\ref{sec:vision-alignment-attention}节）：每个词对每个区域计算一个权重，粒度由参与注意力的单元决定。代价是，越细的粒度越需要相应的监督或自监督信号，例如短语与框的标注、音画同时出现的关系。本章\ref{sec:vision-understanding-retrieval}节在整体粒度上召回候选，再用交叉注意力或逐条件核验处理关系绑定；\ref{sec:vision-understanding-grounding}节分别处理短语与区域、声音与区域、音频与视频片段、指令与控件四种细粒度对应。

\subsection{弱配对与噪声监督带来的学习偏差}
\label{sec:vision-understanding-challenge-weak-pairing}

大规模多模态数据大多不是为当前任务标注的。网页图片的替代文字可能只写“图片1”，教学视频的旁白可能先说“接下来揉面”，几秒后手才开始动；视频的声音可能来自画面之外的人。这类数据只告诉我们“两种模态来自同一来源”，没有说明它们在细粒度上怎样对应，有时甚至并不对应。困难在于：用这样的配对训练，模型学到的可能是与目标无关的共现，例如“厨房背景”与“做饭”一词总是同时出现，模型便只凭背景判断，而不去看手上的动作。

经典做法依靠小规模人工标注：为图像标出类别、为短语标出框，再用固定类别的检测器和人工特征完成匹配。\ref{sec:vision-alignment-weak-pairing}节把这种只在整体层面成立、内部单元对应未知的配对称为弱配对，并区分了假正例与假负例两种噪声；在人工标注的范围内，监督是干净的，但类别和场景的覆盖受标注成本限制。

学习方法用规模和训练目标的设计来对抗噪声。规模大时，单条错配对梯度的影响被大量正确配对冲淡；目标函数可以只要求真实配对比错配高出一个间隔，满足后不再施压，减少对个别噪声样本的过拟合；还可以专门构造困难负例，例如同一视频中的另一句旁白、同一幅图中的背景区域，迫使模型区分细节而非场景。另一类做法是自举：用当前模型给出伪标注，再按可信度加权训练。这些方法降低了标注成本，也带来新的风险：伪标注来自模型自己，错误的高分区域可能被不断强化。本章的HowTo100M（\ref{sec:vision-understanding-howto}节）、声音来源定位（\ref{sec:vision-understanding-sound-source}节）和CRAFT（\ref{sec:vision-understanding-craft}节）分别代表了间隔排序与同源负例、同图困难负例和带置信度的伪标注三种处理方式；VisionTaPas（\ref{sec:vision-understanding-visiontapas}节）用启发式规则生成单元监督，也属于这一类，原论文报告了其中的噪声比例。

\subsection{语言先验带来的捷径与幻觉}
\label{sec:vision-understanding-challenge-language-prior}

问“天空是什么颜色”，不看图回答“蓝色”大多是对的；描述一张街景，写上“路边停着汽车”也常常是对的。语言中的统计规律，即\term{语言先验}（language prior），让模型可以不看视觉证据就给出看似合理的输出。这是一种\term{捷径学习}（shortcut learning）：模型利用了训练数据中与答案相关、却与任务本意无关的线索。它造成两种错误。一种是问题与常识不符时答错，例如阴天的照片仍被答成“蓝色”；另一种是\term{视觉幻觉}（visual hallucination），即输出图中没有依据的对象或属性。OCR也受同样的影响：模糊的“1”和“I”可能被语言上下文改写成常见词，票据编号这样的罕见字符串反而被“纠正”错。

经典流水线中，视觉判断和语言输出是分开的两步，语言部分只能处理视觉部分已经给出的检测结果或字符，先验的作用范围受到限制，但也无法利用语言帮助消歧。\ref{chap:vision-multimodal-alignment}章从融合角度讨论过同一问题：晚期融合（\ref{sec:vision-alignment-late-fusion}节）让各模态分别给出判断再汇总，每个模态的贡献可以单独检查；而模态不平衡（\ref{sec:vision-alignment-imbalance}节）说明，联合训练时更容易拟合的那个模态会主导优化，较弱的模态得不到充分学习。在视觉问答中，问题文字往往就是更容易拟合的一方。

学习方法，尤其是接入大型语言模型之后，语言先验更强了：模型能写出流畅、合乎常识的长段文字，却不一定逐句有视觉依据。应对的方向有三类：在数据上构造反例，例如为同一问题收集答案不同的相似图像；在输出中保留证据，例如让答案附带区域、单元或时间区间；在评价上单独检查，例如询问对象是否存在、替换或移除图像观察答案是否变化。本章\ref{sec:vision-understanding-description}节讨论描述中的幻觉，\ref{sec:vision-understanding-question}节讨论问答中的语言捷径，\ref{sec:vision-understanding-input-protocol}节给出检查模态贡献的方法。

\subsection{高分辨率与长输入带来的证据丢失}
\label{sec:vision-understanding-challenge-resolution}

决定答案的证据常常很小。专业软件的截图上，目标控件平均只占屏幕面积的0.07\%\scite{vision-screenseeker2025}；票据上“128.50”的小数点只有几个像素；一段一小时的视频里，第二次敲门只持续一秒。而视觉编码器的输入大小是固定的，例如224、384或448像素见方，视频也只能抽取有限的帧。把整页、整屏或整段视频压进这个预算，小字、小控件和短暂事件可能在编码之前就已经消失，后面的模型再强也无法恢复。

经典方法对这个问题有成熟的分析。\ref{sec:vision-formation-aliasing}节和\ref{sec:vision-formation-interpolation}节说明，缩小图像时若不先做低通滤波就会产生混叠，而做了低通滤波，小于新采样间隔的细节就被平滑掉了，缩小必然以丢失细节为代价。图像金字塔（\ref{sec:vision-transforms-pyramid}节）在多个尺度上分别处理，经典检测器用多尺度滑动窗口在细尺度上找小目标；视频中，关键帧选择（\ref{sec:vision-videofound-keyframes}节）从冗余的帧中挑出代表帧。这些方法的共同思想是：不在一个尺度上看完全部内容，而是按需要在合适的尺度上看局部。

学习方法沿用了这一思想，并把“看哪里”交给模型或检索来决定。一种做法是改变输入方式，例如按原纵横比切出尽可能多的图像块，而不是把页面压成固定正方形；另一种做法是多次观察，先在全局低分辨率上推测目标所在区域，再裁剪放大；对长视频，则先检索候选片段，或把历史压缩成记忆再按问题回读。代价是计算量随图像块数或观察次数增加，裁剪还可能丢掉判断所需的上下文。本章的ScreenSeekeR（\ref{sec:vision-understanding-screenseeker}节）和Pix2Struct（\ref{sec:vision-understanding-pix2struct}节）分别代表多次局部观察和可变分辨率输入，\ref{sec:vision-understanding-evidence}节讨论页面与长视频中按需获取证据的方法。

\subsection{多步推理带来的证据组合困难}
\label{sec:vision-understanding-challenge-composition}

有些问题不能由一条证据回答。问“甲的2月金额比1月多多少”，要先在表格中找到甲行，再读出两个单元的数值，最后做减法；问“第二次敲门是否在讲话之后”，要先找出所有敲门和讲话的时间区间，再比较第二次敲门的起点与讲话的终点。困难在于，每一步都可能出错，而错误会传递：读错一个数字，后面的计算再准确，答案也是错的；而且只看最终答案，无法知道错误出在哪一步。

经典流水线把每一步写成独立模块，用规则连接：检测图表元素，按刻度把像素换算成数值，再交给确定的程序计算。每一步的输出都可以检查，规则执行的计算也不会出错；但规则只覆盖设计时考虑到的图表样式和问题类型，遇到新样式就失效。

学习方法在“显式步骤”与“端到端”之间给出不同的折中。一种做法让模型只预测结构化的中间决定，例如选择哪些表格单元、执行哪种聚合操作，计算仍由规则完成，中间决定可以检查，但操作集合是预先固定的；另一种做法让模型直接生成答案文字，表达不受限制，却不再暴露中间步骤。本章\ref{sec:vision-understanding-evidence-operations}节讨论怎样用外显的操作组合多条证据，\ref{sec:vision-understanding-answering}节比较从选择、抽取、单元聚合到自由生成四种答案形式在可追溯性与表达范围上的取舍。

表\ref{tab:vision-understanding-challenge-map}把五项挑战与经典方法、学习方法及本章回答的位置对应起来。

\begin{table*}[tbp]
  \centering
  \small
  \begin{tabularx}{\linewidth}{@{}p{24mm}XXX@{}}
    \toprule
    挑战 & 经典方法的处理 & 学习方法的改变 & 本章回答的位置\\
    \midrule
    多粒度对应 & 典型相关分析、动态时间规整、人工特征 & 对比对齐联合学习特征，注意力表达细粒度对应 & 匹配与检索、跨模态定位\\
    弱配对与噪声 & 小规模人工标注、固定类别检测器 & 规模、间隔目标、困难负例、带置信度的伪标注 & HowTo100M、声源定位、CRAFT\\
    语言先验 & 视觉与语言分步处理，晚期融合便于检查 & 更强的语言模型加重先验；反例数据、证据输出与专门评价 & 视觉描述、视觉问答、模态贡献的核验\\
    高分辨率与长输入 & 抗混叠与金字塔、多尺度滑窗、关键帧 & 可变分辨率输入、裁剪放大搜索、检索与记忆 & 界面控件定位、页面解析、证据的获取\\
    证据组合 & 模块化流水线与确定规则 & 单元选择加聚合、自由生成 & 证据的多步组合、答案的形成\\
    \bottomrule
  \end{tabularx}
  \caption{多模态理解的五项挑战与对应处理。经典方法与学习方法并非互斥：后者常常沿用前者的思想，例如裁剪放大搜索延续了多尺度观察，单元聚合延续了规则计算。}
  \label{tab:vision-understanding-challenge-map}
\end{table*}

这些挑战沿着证据从粗到细的顺序依次变得突出，本章的路线也由此确定。检索只需要整体粒度的对应，弱配对的大规模数据在这一粒度上最有用，所以先讨论召回与重排；落到区域、时间和控件时，细粒度对应与分辨率预算成为主要困难，这是定位一节的内容；票据和图表中的证据要先从像素变成字符和结构，才能参与计算，于是读取与解析放在描述和问答之前；描述和问答直接面对语言先验，又需要组合前面各步取得的证据，因此放在后面，并以分层核验收尾。

\section{跨模态匹配与检索}
\label{sec:vision-understanding-retrieval}

\term{跨模态检索}（cross-modal retrieval）指用一种模态的查询，在另一种模态的候选集合中找出相关内容，例如用一句话找图片，或用一句话找视频片段。输入是查询和可能有数十万乃至上亿项的候选库，输出是按分数排列的候选编号。难点在于两个要求相互拉扯：候选库很大，每次查询只能负担很便宜的比较；而“穿红衣服的人正在骑车”这样的请求要求判断对象、属性与关系是否由同一个人满足，这又需要细致的联合计算。

常用的解法是把检索拆成先后三步。第一步规定什么候选算相关，否则模型给出的高分无从验收。第二步召回：查询和候选各自编码成向量，候选向量可以预先计算并缓存，查询到来后只做内积，从全库中取出前$K$个候选。第三步重排：只对这$K$个候选做更昂贵的联合判断，重新排列顺序。召回决定正确结果有没有机会被看到，重排决定它们最终排在哪里；两步的错误要分开检查。

\subsection{检索相关性的规定}
\label{sec:vision-understanding-relevance}

模型给出的分数只是相似程度，是否“找对”要由事先规定的相关性来判断。这一小节要解决的问题是：查询中的哪些条件必须同时成立，候选的单位是什么，正例关系怎样标注。这些规定不依赖具体模型，却决定了召回与重排的评价口径。

“穿红衣服的人正在骑车”同时规定了人物、衣服颜色和动作关系。画面中有红衣行人，也有另一人在骑车，只满足了部分条件；把红色归到自行车上，同样改变了请求。可把完整条件分解为对同一个对象$i$的判断
\[
  \operatorname{person}(i)\land\operatorname{redClothes}(i)
  \land\exists j\,[\operatorname{bicycle}(j)\land\operatorname{ride}(i,j)]\text{。}
\]
这个表达式是相关性协议的说明，不表示双编码器会显式计算这些谓词。它的作用是防止标注和验收把“出现了所有关键词”当成“关系全部成立”。

候选单位也要固定。检索整段视频时，只要片中某处出现目标事件，就可能被标为相关；检索一个两秒片段，则要求该片段包含所需证据。一段文字可能描述多幅图片，一幅图片也可以有多种正确描述，因此应用数据的正例关系通常是多对多。训练划分还应避免同一视频的相邻片段或同一图的近重复版本分散到训练集和测试集。

为诊断组合关系，可以在保持对象词汇的条件下改变请求，如“人搬起箱子”和“箱子挡住人”，或把“戴帽子”改为“没有戴帽子”。相关性应按改变后的条件重新标注。部分相关候选可以保留等级，但不能在一个实验中按完整匹配计分，在另一个实验中又按同主题计分。

\subsection{候选内容的召回}
\label{sec:vision-understanding-candidate-matching}

召回要从$N$个候选中取出前$K$个，$K$通常远小于$N$。为了让每次查询的代价只与一次查询编码和$N$次内积相当，候选必须能够脱离查询单独编码。设候选表示组成矩阵$\bm V\in\mathbb R^{N\times d}$，查询表示为$\bm u\in\mathbb R^{d}$，召回就是取$\bm V\bm u$中最大的$K$项；库很大时还可以用近似最近邻索引代替逐项内积。难点在于：独立编码时查询看不到候选，候选向量必须事先保留所有可能被问到的信息。

双编码召回是\ref{sec:vision-alignment-contrastive}节对比对齐的直接应用。那里的结论是：用批内其他样本作负例、以交叉熵区分真实配对，两个编码器会把配对的表示拉到同一空间中的相近位置，内积因而可以作为相关程度。它的经典前身是\ref{sec:vision-alignment-cca}节的典型相关分析：二者都为两种模态各学一个到公共空间的映射，但典型相关分析作用在固定的人工特征上，映射是线性的，目标是使配对投影的相关系数最大，没有显式的负例；对比对齐则端到端训练深层编码器，并要求真实配对胜过批内其他组合。从融合位置看，两个编码器在计算过程中互不可见，只在最后一步的内积处相遇，相当于放弃了\ref{sec:vision-alignment-intermediate-fusion}节所说的中间交互。这正是双编码能缓存候选向量的原因，也是它难以判断关系绑定的原因。

本小节的三种做法都采用这种双编码结构，区别在于“候选”由什么观测单位构成，以及训练时用什么来源的配对告诉模型哪些内容相关。CLIP以单幅图像为单位，用网页图文对训练；CLIP4Clip把视频看成若干抽样帧，先逐帧编码再沿时间汇总，用人工撰写的视频描述训练；HowTo100M同样以视频片段为单位，但不依赖人工描述，而是把教学视频里自动转写的旁白当作弱配对。三者对同一个难点给出不同回答：图文配对干净但不含时间，人工视频描述含时间但昂贵，旁白规模大但与画面并不总是同步。

\subsubsection{CLIP}
\label{sec:vision-understanding-clip}

CLIP的结构和训练已在\ref{sec:vision-models-contrastive}节展开：图像编码器和文字编码器分别把输入投影到同一个$d$维空间并做$L_2$归一化，训练时在一批真实图文对上计算余弦相似度矩阵，用行、列两个方向的交叉熵把对角线上的配对与批内其他组合区分开。这里用到的结论只有一条：训练后，内积越大表示图文越可能来自同一配对，但这个分数衡量的是批内相对选择，不是对每个条件逐项核对的结果。

用于检索时，先按模型约定裁剪、缩放和归一化图库中的每幅图，编码为$\bm v_i$并存成矩阵$\bm V$。文字查询$q$编码为$\bm u(q)$后，以
\[
  s_i=\bm u(q)^\top\bm v_i
\]
排序并返回图像编号。图像检索文字则反过来缓存候选描述$\bm u_j$，编码查询图像$\bm v$，比较$\bm v^\top\bm u_j$。两种方向共用相似度形式，但候选集合、正例标注和评价分母不同。查询文字可以直接使用用户原句，也可以套用“一张……的照片”这样的模板；模板改变文字编码器看到的上下文，应在比较实验中固定。

假设查询向量为$(1,0)^\top$，三个图像向量为$(0.8,0.6)^\top$、$(0.6,0.8)^\top$和$(0,1)^\top$，均已归一化。教学排序为甲、乙、丙，分数依次为0.8、0.6、0。若人工核对发现乙才满足“红衣人物骑车”的全部关系，甲仅有相关对象，那么模型把部分相关候选排到了完整正例之前。将分数转成softmax并不会补上这个关系判断。这是CLIP召回最典型的失败：整图向量能记住“有红色、有人、有自行车”，却不一定记住红色属于谁。

是否微调取决于目标图库与预训练数据的差距。直接使用原始权重时，不需要目标任务的标注；若用目标领域的图文对微调，通常沿用式\eqref{eq:vision-clip-loss}的双向目标，同时更新两个编码器和投影。此时要注意多对多的正例关系：一个查询的全部已知正例应同时参与目标，或在采样时排除已知正例之间的“假负例”。召回数$K$是主要超参数：$K$越大，后续重排越有机会看到正确候选，重排代价也随$K$线性增加。

\subsubsection{CLIP4Clip}
\label{sec:vision-understanding-clip4clip}

CLIP4Clip的帧编码、时间聚合和训练过程已在\ref{sec:vision-video-clip4clip}节展开。它从CLIP权重出发，把按固定规则抽取的视频帧逐帧编码，再用均值或带时间位置的序列模块汇总成一个视频向量，用视频—描述对上的双向交叉熵微调编码器和聚合器。对检索而言，它把CLIP的“单幅图像”候选换成了“抽样帧的汇总”，因而能比较文字与一段时间内的内容。

使用时，按训练时的帧率和帧数预先编码候选视频或片段，收到查询后与查询向量比较，输出视频编号及分数。抽帧数、片段长度和聚合方式都是需要随结果报告的超参数：抽帧过稀会漏掉短暂动作，片段过长会让目标事件在均值中被稀释。图像已经提供的对象证据与视频增加的时间证据，应分别检查。

对同一图库，依次输入“箱子”“人搬起箱子”“穿红衣服的人搬起箱子”，请求增加了动作和参与者条件。排序改变可能来自真正的时间关系，也可能只来自某帧出现了红衣人物。可以定位命中片段，查看谁接触箱子、箱子是否离开支撑面，以及动作顺序。若帧聚合主要保留整段共同出现的内容，交换“人”和“箱子”的语义角色未必产生足够大的得分变化。

命中整段视频仅完成了候选选择，动作起止仍需时间定位。抽样漏掉搬起瞬间时，下游文字模型可能根据前后场景作合理猜测，但不能把该瞬间写成已经观察到的事实。

\subsubsection{HowTo100M}
\label{sec:vision-understanding-howto}

CLIP4Clip的训练依赖人工为视频片段撰写描述，这类数据的规模受标注成本限制，覆盖的动作种类也有限。教学视频提供了另一种联系：讲解者一边操作一边说明，语音识别生成的字幕行自带起止时间。HowTo100M利用这一点，从约122万段带旁白的教学视频中切出约1.36亿个“片段—旁白”记录，覆盖两万多种任务，不需要人工撰写描述\scite{vision-howto2019}。

这种配对是\emph{弱配对}：旁白可能提前说“接下来揉面”，也可能省略手上正在做的动作，字幕的时间区间与动作并不严格对齐。HowTo100M的主要想法是用规模抵消单条记录的噪声，并在训练中专门加入同一视频内部的错配，让模型必须区分同一厨房里的不同步骤，而不能只凭背景判断“这是做饭视频”。

模型由两个独立的编码分支组成，具体配置见表\ref{tab:vision-understanding-howto-config}。视频一侧使用预先提取、训练时不更新的特征：以ImageNet上训练的ResNet-152每秒取一帧提取外观特征，以Kinetics上训练的16帧ResNeXt-101三维卷积网络每秒约取1.5个特征描述运动；两种特征各自沿片段时间做最大池化，再拼接为一个4096维向量$\bm x_{\mathrm v}$。时间最大池化保留片段内每一维的最强响应，代价是丢掉了这些响应出现的先后顺序。文字一侧以在GoogleNews语料上预训练的word2vec词向量表示旁白中的词，再经一个浅层一维卷积网络得到4096维的$\bm x_{\mathrm t}$；论文没有给出这个卷积网络的层数和卷积核大小。两侧各接一个\term{门控嵌入单元}（gated embedding unit）：它先做线性变换，再让输出的每一维乘上一个由自身计算的0到1之间的系数，按当前内容调节各维的贡献。论文把它写成线性层后接上下文门控。设输入为$\bm x$，有
\begin{equation}
  \bm z=\bm W_1\bm x+\bm b_1,\qquad
  g(\bm x)=\bm z\odot\sigma(\bm W_2\bm z+\bm b_2)\text{，}
  \label{eq:vision-howto-gate}
\end{equation}
其中$\sigma$为逐元素sigmoid，$\odot$为逐元素乘积。原设置中输入维数与输出维数$d$都是4096，因此$\bm W_1,\bm W_2\in\mathbb R^{4096\times4096}$；两侧参数不共享。输出经$L_2$归一化后，以余弦相似度$s_{i,j}$比较第$i$段视频和第$j$条旁白。门控与\ref{sec:vision-alignment-attention}节的注意力都在计算权重，区别在于门控的权重作用在特征维度上，由同一模态的内容决定，并不在两个模态的元素之间建立对应；跨模态的比较仍然只发生在最后的余弦相似度。

\begin{table}[htbp]
  \centering
  \small
  \begin{tabularx}{\linewidth}{@{}p{16mm}XXp{14mm}@{}}
    \toprule
    组件 & 输入 & 结构 & 输出\\
    \midrule
    外观特征 & 每秒1帧 & ImageNet预训练ResNet-152，冻结 & 沿时间最大池化\\
    运动特征 & 16帧片段，每秒约1.5个 & Kinetics预训练ResNeXt-101，冻结 & 沿时间最大池化\\
    视频向量 & 两种特征 & 拼接 & 4096维\\
    文字向量 & 旁白词序列 & word2vec词向量，浅层一维卷积（规模论文未给出） & 4096维\\
    门控嵌入 & 4096维 & 式\eqref{eq:vision-howto-gate}，两侧各一个 & 4096维\\
    \bottomrule
  \end{tabularx}
  \caption{HowTo100M文本—视频嵌入模型的组件配置，依据原论文。两种视觉特征各自的维数论文未分别给出，只给出拼接后的4096维。}
  \label{tab:vision-understanding-howto-config}
\end{table}

训练希望真实配对得分比错配至少高出间隔$\delta>0$。对一个配对$i$及负例$j$，两个方向的损失为
\begin{equation}
  L_{i,j}=[\delta+s_{i,j}-s_{i,i}]_+
          +[\delta+s_{j,i}-s_{i,i}]_+\text{，}
  \qquad [a]_+=\max(0,a)\text{。}
  \label{eq:vision-howto-ranking}
\end{equation}
第一项替换旁白，第二项替换视频。各项已经满足间隔时不再产生惩罚，这与交叉熵始终推动正例概率趋近1不同。原设置取$\delta=0.1$。原实现用批次内的其他记录形成负例，其中一半来自同一段原视频，另一半来自其他视频；否则数量更多的跨视频负例会让模型只学会区分场景。每个批次取32段视频、每段64个片段—旁白对，用学习率$10^{-4}$的Adam训练；预提取的视觉特征保持不变。

以$\delta=0.1$、$s_{i,i}=0.7$、$s_{i,j}=0.65$、$s_{j,i}=0.4$为教学数值，损失为$0.05+0=0.05$。若$j$是同一厨房里的另一条旁白，第一项要求模型把两个步骤分开，只认出厨房背景还不够。

训练后的两个映射可以直接用于新数据：对查询编码，按片段分数排序，再把片段时间还原到原视频。原论文既在YouCook2、CrossTask等教学视频数据上不经微调直接检索和定位步骤，也在MSR-VTT等人工描述数据上以HowTo100M预训练为起点再微调。排序目标鼓励区别配对，却不会自动消除旁白和实际动作之间的错位；\ref{sec:vision-alignment-contrastive}节介绍的MIL-NCE就是针对这一点的后续改进，它把时间上相邻的若干段旁白都当作候选正例；查询内容若是教学视频很少涉及的场景，例如街头监控，弱配对数据的覆盖也帮不上忙。

三种召回方式的差别可以放在同一张表中对照（表\ref{tab:vision-understanding-recall-compare}）。从CLIP到CLIP4Clip，改变的是候选单位：把单幅图换成多帧汇总，换来了对时间内容的比较能力，代价是抽帧与聚合会稀释短暂事件。从CLIP4Clip到HowTo100M，改变的是配对来源：用自动旁白代替人工描述，换来了规模，代价是配对噪声和领域局限。三者都只能给出整体相似度，关系是否由同一对象满足，要留给重排处理。

\begin{table}[htbp]
  \centering
  \small
  \begin{tabularx}{\linewidth}{@{}p{14mm}XXX@{}}
    \toprule
    比较项 & CLIP & CLIP4Clip & HowTo100M\\
    \midrule
    候选单位 & 单幅图像 & 抽样帧汇总的片段 & 字幕时间区间对应的片段\\
    配对来源 & 网页图文对 & 人工视频描述 & 自动转写的旁白\\
    视觉编码 & 图像编码器 & CLIP图像编码器逐帧编码后聚合 & 预提取的二维与三维卷积特征\\
    训练目标 & 双向交叉熵 & 双向交叉熵 & 双向间隔排序\\
    典型失败 & 属性与对象错绑 & 短暂动作被稀释 & 旁白与动作错位\\
    \bottomrule
  \end{tabularx}
  \caption{三种召回方式的对照。三者都预先编码候选、以内积排序；表中各行说明它们在候选单位、配对监督和失败情形上的差别。}
  \label{tab:vision-understanding-recall-compare}
\end{table}

\subsection{候选内容的重排}
\label{sec:vision-understanding-reranking}

召回给出$K$个候选，其中可能有同主题而关系不同的内容，例如红衣行人和蓝衣骑车者同框的图片。重排的输入是查询和这$K$个候选，输出是新的顺序，难点是让查询直接作用于候选的局部内容，判断条件是否由同一对象满足。因为只处理$K$项，可以负担每项一次完整的联合计算。

按\ref{chap:vision-multimodal-alignment}章的融合位置来看，重排是在召回放弃的地方补上中间融合（\ref{sec:vision-alignment-intermediate-fusion}节）：查询的词元与候选的视觉单元在网络内部通过交叉注意力相互读取，每个词可以把注意力放到它所描述的区域上，关系绑定才有机会在计算中体现。中间融合的代价是表示不能离线缓存，每个查询—候选对都要重新计算一次，所以只能用于召回之后的少量候选。

重排有两条路线，区别在于关系判断是否由一个学到的分数隐式完成。基于配对判定的重排把查询和候选一起送入交叉注意力网络，输出一个“是否匹配”的概率，BLIP-2的图文匹配头是这种做法。基于定位与读取的重排把查询拆成可检查的条件，调用定位和读取接口取得区域证据，再按条件满足情况排序。前者一次计算、无需拆解查询，后者每个条件都有可核对的依据，但依赖拆解规则和下游接口。两条路线都无法找回召回阶段已经漏掉的正例。

\subsubsection{基于配对判定的重排}
\label{sec:vision-understanding-blip-retrieval}

BLIP-2的Q-Former已在\ref{sec:vision-models-blip2}节展开：32个可学习查询通过交叉注意力读取图像编码器输出，形成固定数量的图像查询表示；第一阶段预训练用同一组参数、三种注意力掩码学习三个目标。其中图文对比目标让查询与文字互不可见，得到可以独立缓存的表示；图文匹配目标让查询与文字双向交互，由二分类头判断图文是否配对。检索正好用到这两种模式：前者负责召回，后者负责重排。

使用时，先用图像查询表示与文字表示的相似度召回有限候选。重排阶段把一个候选文字与图像查询一同送入Q-Former，并采用允许二者互看的注意力掩码；每个查询输出经二分类头得到匹配logit，取平均后得到配对概率$p_1$，按$p_1$重新排列候选。这条路径只用到Q-Former，不需要大型语言模型逐词生成。由于每个候选都需要一次联合前向计算，文字检索图像时代价是$K$次Q-Former计算，而召回只需一次查询编码。

原论文的检索适配共同微调图像编码器和Q-Former，并保留对比、匹配及图像条件文字生成三个目标；这不同于第一阶段冻结图像编码器的预训练。其推理设置先取128个候选，再按配对匹配分数重排。冻结范围和候选数都应随所用实验设置报告，不能从预训练说明中直接沿用。

重排能够让“同一画面里的红衣人物是否就是骑车者”影响联合计算，但仍需关系数据检验：匹配头的训练负例多数来自不同图片，未必覆盖属性在同图内换绑的情形。若图库有三个完整正例，召回集仅含其中两个，那么无论重排多准确，完整正例的召回率最多为$2/3$。这个上限来自候选已经被删除，与后续模型有多少参数无关。

\subsubsection{基于定位与读取的重排}
\label{sec:vision-understanding-constraint-check}

配对概率是一个数，它不说明哪个条件成立、哪个不成立。对需要严格满足条件的查询，可以使用一个教学组合流程：从请求中列出需要核对的对象、属性、关系、计数或文字条件$c_1,\ldots,c_m$；对高分候选调用后文的定位和读取接口取得区域证据；再对每个条件记录状态$z_k\in\{\text{满足},\text{不满足},\text{证据不足}\}$。排序时先比较“不满足”的条件数，少者在前；再比较“证据不足”的条件数；仍相同时才用召回分数。这条规则不新增模型训练，也不属于CLIP或BLIP-2的原始输出。

例如，甲图得分0.8，但定位后发现红衣人物在步行，骑车者穿蓝衣，记录为“动作绑定不满足”，排到后面。乙图得分0.6，相关人物和自行车位置、接触关系均可见，三项条件都满足，排在前面。若车身遮住腿部，图像又无法区分骑行和推行，应把这一条记为证据不足，而不是用较高相似度补足不可见部分。

这一流程的超参数是参与核验的候选数和各接口的阈值，失败主要来自两处：拆解规则漏掉了条件，或者定位与读取接口本身出错。结果可以同时返回候选编号、匹配分数、目标区域及各条件状态，使排序错误能被追溯到具体环节。

两种重排可以按同一查询和同一召回集比较（表\ref{tab:vision-understanding-rerank-compare}）。基于定位与读取的重排把配对判定的单个分数拆成逐条件的证据记录，换来了可解释性和对“证据不足”的显式表达，代价是更多的接口调用，以及对查询拆解规则的依赖；配对判定不需要拆解，但它犯错时只能从数据和模型训练上找原因。

\begin{table}[htbp]
  \centering
  \small
  \begin{tabularx}{\linewidth}{@{}p{16mm}XX@{}}
    \toprule
    比较项 & 基于配对判定的重排 & 基于定位与读取的重排\\
    \midrule
    每个候选的计算 & 一次联合前向计算 & 每个条件一次定位或读取\\
    输出 & 配对概率 & 各条件的状态与区域证据\\
    训练需求 & 需要图文匹配训练 & 复用已有接口，不新增训练\\
    能否表达证据不足 & 不能，只有一个分数 & 能，单独记录\\
    主要失败 & 匹配头未见过的换绑情形 & 拆解遗漏条件、接口误差\\
    \bottomrule
  \end{tabularx}
  \caption{两种重排的对照。两者都在同一召回集上工作，都不能找回召回阶段漏掉的正例。}
  \label{tab:vision-understanding-rerank-compare}
\end{table}

\section{跨模态定位}
\label{sec:vision-understanding-grounding}
\label{sec:vision-understanding-grounding-query}

检索给出“哪幅图相关”，而“右边拿杯子的人”需要落到画面中的某个对象。\term{视觉定位}（visual grounding）把文字或其他查询对应到视觉对象、空间区域或时间位置。输入是一份观测和一个查询，输出是框、掩码、点或时间区间；输出的位置必须仍属于原始观测，才能用于复核、裁剪或后续操作。

定位之前要先解析查询，确定目标是什么、靠什么参照来区分。在“右边拿杯子的人”中，“人”是目标类别，“拿杯子”是人与杯子的关系，“右边”还依赖参照系。若两个人都在拿杯子，“右边”可能指图像右侧，也可能指说话者右手方向。系统应把默认参照系写入协议；本章未另说明时，图像横轴向右、纵轴向下。多图输入还需指定“图2中的人物”，页面输入需带页码，视频输入需带时刻或区间。“他旁边的杯子”依赖此前已经确定的“他”，不能仅按最近出现的名词替换。无法唯一确定时，可以返回候选集合或请求补充参照条件；查询目标根本不存在，也应允许输出空结果。

输出形式也随任务规定。框表示外接范围，掩码描述像素归属，点击点必须位于可操作控件内，区间描述事件持续时间。它们的几何关系可以辅助转换，但不能仅凭输出了四个数就认定完成框定位，也不能把框中心总是可点击当作普遍条件。

查询与被定位的对象可以属于不同模态，本节按定位对象展开四种情形。文字短语定位到图像区域时，监督最充分，可以端到端训练检测式模型。声音来源定位到图像区域时，没有框标注，只能从音画同时出现的关系中自监督学习。音画事件定位到时间区间时，要比较两条时间序列。界面控件按指令定位到点击位置时，截图分辨率高、目标极小，难点转向观测的分辨率预算。四种情形的共同问题是：定位需要保留空间或时间分辨率，把整图压成一个向量再与整句比较，无法知道哪个位置支持哪段表达。

\subsection{文字短语的区域定位}
\label{sec:vision-understanding-grounding-match}

输入是一幅图像和一段文字，文字可以是类别名、短语，也可以是\term{指代表达}（referring expression），即用属性和关系唯一指明某个对象的描述，如“右侧戴帽子的人”。输出是若干框，每个框附带它对应的文字位置。难点有两个：文字词汇是开放的，不能事先列出类别表；同一句话可能提到多个对象，模型要知道每个框对应句中哪几个词。

两种代表模型都沿用DETR式集合预测（\ref{sec:vision-recognition-detr}节）：一组对象查询经解码得到框，训练时用二分匹配把预测与标注一一对应。它们的差别在于文字在何处、以何种方式进入检测计算。MDETR在编码器入口把图像特征与文字词元拼成一个序列，此后每一层都联合计算，并让每个对象查询预测它指向哪些词元位置。Grounding DINO保留两条分支，在特征增强阶段交换信息，再用文字相关性挑选初始化对象查询，并以更大规模的检测与定位数据训练，面向开放类别检测。

用\ref{chap:vision-multimodal-alignment}章的术语说，两者都属于中间融合（\ref{sec:vision-alignment-intermediate-fusion}节）：各模态先由自己的编码器处理，再在网络内部交互。MDETR的交互方式是把两种特征拼接成一个序列后做自注意力，形式上接近\ref{sec:vision-alignment-early-fusion}节的特征拼接，但拼接发生在各自编码之后，而且拼接后还有多层注意力，每个词元都能对每个视觉单元分配权重；Grounding DINO则用两个方向的交叉注意力交换信息，两条分支始终保持各自的表示。这两种结构都让词与区域之间的注意力权重成为\ref{sec:vision-alignment-attention}节所说的软对齐，而短语—框标注直接监督了这种对齐的结果。

\subsubsection{MDETR}
\label{sec:vision-understanding-mdetr}

在MDETR之前，常见的文字定位系统先用一个固定类别的检测器提出候选框，再让文字与这些框匹配。检测器若没有为“行李架”这类不在类别表中的对象提出候选，后续匹配无从补救；检测器训练时也不知道当前句子要找什么。MDETR的主要想法是让文字从一开始就调制检测：图像与文字在同一个Transformer编码器中交换信息，对象查询直接读取融合后的表示产生框，因而候选框的生成本身就以文字为条件\scite{vision-mdetr2021}。

模型的数据流如下。图像经卷积骨干得到$h\times w$的空间特征网格，用$1\times1$卷积降到$d$维后展平为$hw$个视觉单元，并加上二维位置编码以保留空间信息。文字经预训练的RoBERTa编码为$l$个词元表示，再线性投影到同样的$d$维。两组单元沿序列方向拼接成$(hw+l)\times d$的矩阵，送入联合Transformer编码器，自注意力在视觉单元和词元之间不加区分地交换信息。解码器有$N$个可学习的对象查询，读取编码器输出后，每个查询产生一个$d$维表示$\bm o_i$，再接两个预测头：框头输出归一化的中心与宽高，词元头输出一个在全部词元位置及额外“无对象”位置上的概率分布$p_{i,\cdot}$。词元位置的上限为$L=256$。

表\ref{tab:vision-understanding-mdetr-config}列出各组件的规模。论文正文给出了文字编码器、骨干、查询数、温度与训练轮数；联合编码器与解码器的层数、宽度和头数，正文只说沿用DETR，具体数值取自官方实现的默认配置，表中分别注明。以ResNet-101为例作一个教学计算：骨干最后一级特征图的步长为32，若输入图像为$800\times1333$像素，特征图约为$25\times42$，即约1050个视觉单元；加上至多256个词元，联合编码器处理的序列长度约为1300。论文正文未给出预训练时的图像尺寸，这里的尺寸只用于说明序列长度的量级。

\begin{table}[htbp]
  \centering
  \small
  \begin{tabularx}{\linewidth}{@{}p{18mm}XX@{}}
    \toprule
    组件 & 配置 & 来源\\
    \midrule
    图像骨干 & ResNet-101（ImageNet预训练，冻结批归一化），或EfficientNet-B3（约1200万参数）、B5（约3000万参数） & 论文\\
    文字编码器 & RoBERTa-base：12层，宽度768，12头 & 论文\\
    投影与位置 & 两侧线性投影到$d=256$；图像加二维位置编码 & 投影见论文，$d$见官方实现\\
    联合编码器 & 6层，宽度256，8头，前馈维数2048，丢弃率0.1 & 官方实现\\
    解码器 & 6层，100个对象查询，宽度与头数同编码器 & 查询数见论文，层数见官方实现\\
    预测头 & 框：中心与宽高；词元：256个位置加“无对象” & 论文\\
    对比对齐 & 投影到64维，温度$\tau=0.07$ & 维数见官方实现，温度见论文\\
    预训练 & 约130万对图文、约20万幅图像；40轮，等效批大小64 & 论文\\
    参数量 & 论文未给出 & —\\
    \bottomrule
  \end{tabularx}
  \caption{MDETR的组件配置。官方实现指作者发布代码中的默认参数；论文正文未给出的数值不作推测。}
  \label{tab:vision-understanding-mdetr-config}
\end{table}

联合编码器沿用DETR的Transformer块，默认在残差相加之后做层归一化。设输入序列为$\bm Z\in\mathbb R^{(hw+l)\times256}$，位置编码为$\bm P$，单个编码器块计算
\[
  \bm Z'=\operatorname{LN}\bigl(\bm Z+\operatorname{MSA}(\bm Z+\bm P,\,\bm Z+\bm P,\,\bm Z)\bigr),\qquad
  \bm Z''=\operatorname{LN}\bigl(\bm Z'+\operatorname{FFN}(\bm Z')\bigr)\text{，}
\]
其中$\operatorname{MSA}(\cdot,\cdot,\cdot)$的三个参数依次为查询、键和值，位置编码只加在查询和键上；$\operatorname{FFN}$为$256\to2048\to256$的两层全连接网络。解码器块在自注意力和前馈网络之间多一层交叉注意力，以对象查询为查询、以编码器输出为键和值。对象查询读到的编码器输出已经混合了文字信息，所以框的产生本身是以文字为条件的。

训练数据是带有“短语—框”对应的图文记录，原论文从Flickr30k、COCO的指代表达数据和Visual Genome中汇集约130万对图文。集合匹配沿用DETR，在预测框与真实框之间寻找一对一分配，再计算框的$L_1$损失和广义IoU损失。文字对应则允许多对多：一个短语可以指两个人，同一个人也可能在句中被不同短语指代。

第一个文字目标是\term{软词元预测}（soft token prediction）：不要求模型预测类别名，而是预测当前对象在句中对应的词元位置分布。设一个已匹配对象对应的词元位置集合为$T_i^+$，目标把概率质量均匀分给这些位置，其交叉熵为
\begin{equation}
  L_{\mathrm{token},i}
  =-\frac{1}{|T_i^+|}\sum_{j\in T_i^+}\log p_{i,j}\text{。}
  \label{eq:vision-mdetr-token}
\end{equation}
例如“红衣／人物”占两个词元位置，目标分别为$1/2$，不是要求模型任选一个位置得到概率1。未匹配的对象查询则把概率集中到“无对象”位置。实际对应建立在分词后的位置上，不能用中文字符数代替模型词元数。

软词元预测约束的是输出位置，并没有直接约束表示之间的距离。MDETR另外把对象表示$\bm o_i$与编码器输出的词元表示分别投影到低维并归一化，得到$\hat{\bm o}_i$与$\hat{\bm t}_j$，计算对象到文字的对比项
\begin{equation}
  L_{\mathrm{o\to t},i}
  =-\frac{1}{|T_i^+|}\sum_{j\in T_i^+}
    \log\frac{\exp(\hat{\bm o}_i^\top\hat{\bm t}_j/\tau)}
    {\sum_k\exp(\hat{\bm o}_i^\top\hat{\bm t}_k/\tau)}\text{。}
  \label{eq:vision-mdetr-contrast}
\end{equation}
文字到对象方向按相同方式对齐，再取两个方向的平均；没有正对应的项不参与构造。原论文按文献惯例取温度$\tau=0.07$。总目标为框的$L_1$与广义IoU损失、软词元损失和对比对齐损失的加权和，官方实现默认的权重依次为5、2、1、1，共同更新骨干、文字编码器、联合编码器和解码器；预训练时骨干与Transformer的学习率按论文附录分别为$10^{-4}$与$10^{-5}$，35轮后降为十分之一；文字编码器的学习率先线性预热到$5\times10^{-5}$，再线性衰减。两个文字目标使用相同标注，但作用不同：一个规定“输出指向哪里”，一个规定“表示彼此靠近”。

预训练后，MDETR可以直接用于短语定位：对句中每个目标短语，把各对象查询在该短语词元上的概率相加，按这个分数筛选框并恢复原图坐标。用于指代表达理解、GQA式问答等下游任务时，原论文在相应数据上继续微调。若图中有红衣人物甲和蓝衣人物乙，输入“红衣人物拿着杯子，蓝衣人物站在旁边”，模型可以同时输出两个人物框及对应短语位置；若任务只要求拿杯者，还需排除作为参照出现的乙。框能支持区域检查，像素掩码需要另加分割路径；复杂关系能否正确消歧，仍取决于训练与测试数据是否覆盖相应关系。

\subsubsection{Grounding DINO}
\label{sec:vision-understanding-grounding-dino}

Grounding DINO的跨模态特征增强、语言引导的查询选择和跨模态解码已经在\ref{sec:vision-recognition-grounding-dino}节展开。与MDETR相比，它保留图像与文字两条分支，在多个阶段做双向交叉注意力，并按候选位置与词元的相似度挑选初始化对象查询；输出同样是框和框对各词元的匹配得分。这里讨论把它用于文字短语定位时，输入怎样组织、输出怎样解码。

输入文字有两种组织方式。类别词列表按所用接口的分隔约定写成“人物 . 帽子 .”这样的串，每个类别是一段独立的短语；完整指代表达则作为一句话输入，保留属性和关系，不随意拆成互不相关的词。把“人物”“戴帽子的人”“右侧戴帽子的人”分别作为查询，可以观察属性和参照条件对结果的影响。使用已发布检查点时，通常不为每次查询更新权重，并保留其分词、图像尺度和阈值协议；原论文在指代表达任务上另做了微调，两种用法的结果不应混为一谈。

解码时，先取每个框在全部词元上的最高得分，以框阈值筛出候选；再对每个保留的框，以文字阈值选出得分足够高的词元，拼接恢复短语。两道阈值作用不同：降低框阈值会增加候选区域，降低文字阈值会扩大附着在框上的文字范围。二者都不能把“右侧”一定指向当前目标作为逻辑保证，返回的短语也未必是连续、完整的语法单位。对“右侧戴帽子的人”，应核对帽子属于哪个人物、“右侧”的参照是否一致，以及是否存在第二个满足条件的候选。需要概念对应的实例掩码时，可以使用\ref{sec:vision-recognition-sam3}节的SAM 3接口，分别检查概念存在与实例范围；框和掩码的组合增加了成果类型，并没有免除关系核验。

两者在同一任务上的差别集中在文字进入检测的位置。MDETR从编码器入口开始联合计算，词元与视觉单元在每一层都能直接交互，适合需要逐词对应的短语定位，但每次查询都要重新运行整个编码器，训练数据也局限于带短语对应的图文集。Grounding DINO把文字作用分散到特征增强、查询选择和解码三处，并利用大规模检测数据，换来更强的开放类别检测能力；代价是长句中的关系仍主要靠词元匹配得分表达，阈值解码可能截断短语。两者都输出框而非掩码，也都需要另外检查参照关系。

\subsection{声音来源的区域定位}
\label{sec:vision-understanding-sound-source}

一帧画面中同时出现说话者和旁观者，问题是“声音来自画面中哪里”。输入是一帧图像和与之同时的一段音频，输出是一张与图像网格对齐的热图，高值区域表示\term{声源定位}（sound source localization）的结果。难点在于几乎没有带框的声源标注，能大量获得的只有“这段声音与这段画面来自同一视频”的配对关系。因此这一任务通常采用自监督路线：让模型从音画是否同时出现中学习，再从学到的空间相似度中读出位置。路线之间的区别在于负例怎样构造。

经典方法把这个问题看成像素与声音之间的相关分析。发声物体的运动与声音的强弱同步变化，例如说话者的嘴唇开合与语音能量起伏一致，因此可以在一段视频上，统计每个像素的亮度变化与音频特征随时间的相关程度，相关最强的位置就是声源。Pixels that Sound把\ref{sec:vision-alignment-cca}节的典型相关分析用于这一问题，并要求视觉一侧的投影稀疏，使高权重集中在少数像素上\scite{vision-pixelssound2005}。这类方法不需要标注，但依赖声音与像素变化的同步，对静止的发声物体或背景中同样在动的物体都会失效。学习方法把像素换成深层特征，把时间相关换成音画是否同时出现的对比判别，从大量视频中学到“什么样子的东西发出什么样的声音”，因而能定位单帧中静止的发声物体。

此前常见的学习做法把整幅图与声音作对比：同一视频的图和声为正例，不同视频的为负例。这种训练只要求区分“有人说话的画面”和“另一种场景”，并不要求模型指出画面中哪个位置在发声，定位能力只是附带出现。Localizing Visual Sounds the Hard Way的主要想法是把同一幅图中的背景区域也当作负例：既然画面里只有一部分在发声，模型就应让声音与这部分的相似度高于与同图其他部分的相似度\scite{vision-visual-sounds2021}。

模型有两条编码分支。图像经卷积网络保留空间网格，得到各位置的特征$\bm f_{u,v}\in\mathbb R^d$；音频先转为频谱图，再经另一卷积网络并全局池化为单个向量$\bm a\in\mathbb R^d$。对非零的$\bm a$与$\bm f_{u,v}$，逐位置计算余弦相似度
\begin{equation}
  S_{u,v}=\frac{\bm a^\top\bm f_{u,v}}
                 {\lVert\bm a\rVert_2\lVert\bm f_{u,v}\rVert_2}\text{，}
  \label{eq:vision-sound-map}
\end{equation}
得到一张$h\times w$的热图。声音编码器输出整段表示，因此热图描述的是这段声音与当前图像的空间联系，不是每个音频采样点的位置。这个计算与\ref{sec:vision-alignment-attention}节的注意力打分形式相同：音频向量相当于查询，各空间位置的视觉特征相当于键，只是这里不做softmax归一化，而直接把余弦相似度当作定位分数。

表\ref{tab:vision-understanding-lvs-config}给出原论文的配置。两条分支都用ResNet-18：图像分支保留空间网格，输入$224\times224\times3$时输出$14\times14\times512$的特征；音频分支把3秒音频转成$257\times300$的幅度谱图，经ResNet-18得到$17\times13\times512$的特征图，再做最大池化得到512维向量。由于两侧通道数相同，式\eqref{eq:vision-sound-map}可以直接在512维上计算。

\begin{table}[htbp]
  \centering
  \small
  \begin{tabularx}{\linewidth}{@{}p{14mm}XXX@{}}
    \toprule
    分支 & 输入 & 主干 & 输出\\
    \midrule
    图像 & $224\times224\times3$，不裁剪 & ResNet-18，ImageNet预训练 & $14\times14\times512$特征网格\\
    音频 & 3秒片段的$257\times300$幅度谱图 & ResNet-18 & $17\times13\times512$，最大池化为512维\\
    热图 & 两侧输出 & 逐位置余弦相似度 & $14\times14$\\
    \bottomrule
  \end{tabularx}
  \caption{Localizing Visual Sounds the Hard Way的结构配置。论文未给出音频采样率；测试时使用完整长度的谱图。}
  \label{tab:vision-understanding-lvs-config}
\end{table}

训练只有音画配对时，真实声源掩码不可用。方法用当前热图本身构造较可信的正区域和同图负区域，用两道阈值$\epsilon_+>\epsilon_-$和温度$\tau$得到软掩码：
\[
  m^+_{u,v}=\sigma\!\left(\frac{S_{u,v}-\epsilon_+}{\tau}\right),
  \qquad
  m^-_{u,v}=1-\sigma\!\left(\frac{S_{u,v}-\epsilon_-}{\tau}\right)\text{。}
\]
高分位置获得较大正权重，低分位置获得较大负权重，中间区域对两侧的贡献较小；由于使用sigmoid，中间带不是严格为零。对批次中的第$i$个视频，正响应$P_i$是以$m^+$加权的平均相似度；负响应$N_i$由两部分相加：以$m^-$加权的同图背景平均相似度，以及该声音与批内其他视频图像各位置相似度之和除以网格大小$hw$。训练目标为
\begin{equation}
  L=-\frac1k\sum_{i=1}^{k}\log\frac{\exp(P_i)}{\exp(P_i)+\exp(N_i)}\text{，}
  \label{eq:vision-sound-loss}
\end{equation}
其中$k$为批次大小。它提高正响应相对于两类负响应的得分，同时更新图像和音频编码器。原设置取$\epsilon_+=0.65$、$\epsilon_-=0.4$、$\tau=0.03$：温度很小，sigmoid接近阶跃函数，相似度高于0.65的位置几乎全部算作正区域，低于0.4的几乎全部算作负区域。训练使用学习率$10^{-4}$的Adam，批大小256，数据取自Flickr SoundNet和VGGSound中1万或14.4万对图像与音频；原论文还为评价构建了VGG-SS测试集，表中列出5158段视频、220个类别，正文概称约五千段、两百多类。

使用时，把热图上采样到原图尺寸，按所定阈值得到区域，再与标注的可见声源区域比较。同图困难负例使训练必须区分发声者和视觉上同样显著的旁观者。不过，正负区域来自模型自己的热图，这种自举仍可能把错误的高分区域当成正例并不断强化。若声音来自画外，热图峰值仍会指向某个可见对象；处理这种情况需要另设“无可见声源”判定和相应数据，原热图接口不能给出保证。

\subsection{音画事件的时间定位}
\label{sec:vision-understanding-av-time}

空间声源回答“画面哪里”，时间定位回答“另一模态的哪一段与这段内容同步”。输入是一段音频查询和一段较长的视频（或反过来），输出是视频中与之对应的时间区间。难点在于两条序列的单元要能直接比较，而训练时只知道哪些音画片段同时发生。Audio-Visual Event Localization一文为此构建了AVE数据集，其中的视频包含同时可见又可听的事件，并提出从音频定位视觉片段及其反方向的跨模态定位任务\scite{vision-ave2018}。

该文的主要想法是学习一个音画距离：同步的音频片段与视觉片段在共同空间中距离小，错配的距离大，定位时沿时间滑动寻找累计距离最小的窗口。AVE数据集含4143段10秒视频，覆盖28类事件；每段视频按1秒切成10个单元，每个单元标注事件类别或背景。

模型称为音画距离学习网络，结构见表\ref{tab:vision-understanding-avdln-config}。视觉一侧，每个1秒单元取16帧RGB图像，用ImageNet上预训练的VGG-19提取pool5特征图，再在16帧上取平均，得到一个$7\times7\times512$的特征图；音频一侧，用在AudioSet上预训练的类VGG音频网络为每个1秒单元提取128维表示。两种特征都不再更新。两侧各接一个两层全连接网络，投影到共同空间后以欧氏距离$D=\lVert\bm r_{\mathrm v}-\bm r_{\mathrm a}\rVert_2$比较，其中$\bm r_{\mathrm v}$、$\bm r_{\mathrm a}$为两侧的投影结果。论文没有给出全连接层的宽度，也没有说明$7\times7\times512$的视觉特征在进入全连接层之前怎样展平或池化。

\begin{table}[htbp]
  \centering
  \small
  \begin{tabularx}{\linewidth}{@{}p{14mm}XXX@{}}
    \toprule
    分支 & 输入 & 特征提取 & 投影\\
    \midrule
    视觉 & 每秒16帧RGB & ImageNet预训练VGG-19的pool5，16帧平均为$7\times7\times512$ & 两层全连接，宽度论文未给出\\
    音频 & 每秒一段音频 & AudioSet预训练的类VGG网络，128维 & 两层全连接，宽度论文未给出\\
    距离 & 两侧投影 & — & 欧氏距离，对比损失间隔$m=2.0$\\
    \bottomrule
  \end{tabularx}
  \caption{AVE音画距离学习网络的结构配置，依据原论文。两种特征提取网络均为预训练后固定的模块。}
  \label{tab:vision-understanding-avdln-config}
\end{table}

若$y=1$表示同步配对，$y=0$表示错配，训练目标为
\begin{equation}
  L=yD^2+(1-y)[m-D]_+^2\text{，}\qquad m>0\text{。}
  \label{eq:vision-av-distance}
\end{equation}
同步对被拉近，错配对在距离小于$m$时被推开，超过$m$后不再受罚；原设置取$m=2.0$。标签来自训练配对关系，不是模型使用时预先知道目标时间。

设音频查询含$l$个时间单元，候选视频含$T$个时间单元。沿视频滑动同样长度的窗口，选择累计距离最小的起点
\begin{equation}
  \hat r=\argmin_{1\leq r\leq T-l+1}
    \sum_{j=1}^{l}D(\bm v_{r+j-1},\bm a_j)\text{。}
  \label{eq:vision-av-sliding}
\end{equation}
这里$\bm v_k$、$\bm a_j$表示已投影的片段特征。若单元时长为$\Delta$，且视频从0秒开始，第$r$个单元覆盖$[(r-1)\Delta,r\Delta)$，输出区间为$[(\hat r-1)\Delta,(\hat r+l-1)\Delta)$。

图\ref{fig:vision-av-temporal-match}用二维投影的教学数据走通三次窗口比较。取音频向量$(0,0)$、$(1,0)$，四段视频向量$(2,0)$、$(0,0.1)$、$(1,0.2)$、$(3,0)$。三个窗口的累计距离约为3.005、0.300和3.020，故选中第2、3单元；当$\Delta=1$秒时，区间为$[1,3)$秒。

\begin{figure}[htbp]
  \centering
  % 教学欧氏距离：声音(0,0),(1,0)，视频(2,0),(0,.1),(1,.2),(3,0)。
\begin{tikzpicture}[x=1cm,y=1cm,font=\figurefont\small]
  \node[anchor=west] at (0,3.5) {逐单元距离$D(\bm v_k,\bm a_j)$};
  \foreach \x/\k in {2.2/1,4.2/2,6.2/3,8.2/4}
    \node at (\x,2.7) {视频$\k$};
  \node[anchor=east] at (1.1,1.8) {声音1};
  \node[anchor=east] at (1.1,0.7) {声音2};
  \foreach \x/\y/\v in {
    2.2/1.8/2.000,4.2/1.8/0.100,6.2/1.8/1.020,8.2/1.8/3.000,
    2.2/0.7/1.000,4.2/0.7/1.005,6.2/0.7/0.200,8.2/0.7/2.000}{
    \draw[BookRule,line width=0.5pt,fill=BookPaper]
      (\x-0.9,\y-0.4) rectangle (\x+0.9,\y+0.4);
    \node at (\x,\y) {$\v$};
  }
  \draw[BookTeal,line width=1.2pt] (3.3,1.4) rectangle (5.1,2.2);
  \draw[BookTeal,line width=1.2pt] (5.3,0.3) rectangle (7.1,1.1);
  \node[anchor=west] at (0,-0.2) {窗口累计距离：};
  \node[align=center] at (1.7,-1.0) {视频1--2\\$3.005$};
  \node[align=center,text=BookTeal] at (4.8,-1.0) {视频2--3\\$\mathbf{0.300}$};
  \node[align=center] at (7.9,-1.0) {视频3--4\\$3.020$};
  \node[anchor=west] at (0,-2.0) {选中$[1,3)$秒；时间单元长1秒};
\end{tikzpicture}

  \caption{音频查询与视频窗口对应的教学计算。每列是一个视频时间单元；窗口内按相对次序逐项相加，粗框标出累计距离最小的窗口。向量和距离均为构造值，不是模型测量结果。}
  \label{fig:vision-av-temporal-match}
\end{figure}
\FloatBarrier

式\eqref{eq:vision-av-sliding}可以看成\ref{sec:vision-alignment-dtw}节动态时间规整的一个受限特例。动态时间规整允许一条序列的一个单元对应另一条序列的多个单元，用动态规划在所有单调路径中找累计代价最小的一条，适合两条序列速度不同的情形，例如同一句话说得快慢不同。滑动窗口只允许对角线路径：音频的第$j$个单元固定对应窗口内的第$j$个视频单元，于是只需比较$T-l+1$个起点。这个限制在这里是合理的，因为音频与视频来自同一时钟，事件的持续时间在两种模态中相同；若查询来自另一段录音，节奏与目标视频不同，就需要放开对应路径。原论文的评价只在事件短于10秒的测试视频上进行，并要求预测的片段与真实片段完全一致才算正确。

重复敲门可能在两处形成同样小的距离，最小值不再唯一；较粗的单元又限制边界精度。原工作另外研究逐秒预测“同时可见且可听的事件类别”，其监督是事件标签，输出是时间标签序列。该任务与上述查询对应的监督及输出均不同，不能把两套结果相互替代。

\subsection{界面控件的指令定位}
\label{sec:vision-understanding-context}

在界面上执行“保存当前文档”，需要把指令落到截图中的某个控件。\term{界面定位}（GUI grounding）的输入是一张截图和一条自然语言指令，输出是点击点或控件框。它与自然图像定位的难点不同：截图常有上千像素宽，目标控件可能只占极小面积；界面上文字密集，可能有两个“保存”按钮，一个属于当前对话框，另一个属于背景窗口。整屏缩小后，决定控件区别的小字可能消失。

两种代表做法分别从模型和观测两端解决这个问题。SeeClick通过大规模界面定位数据训练视觉语言模型，让它一次读完整张截图就直接生成坐标；它改进的是模型本身的定位能力。ScreenSeekeR不训练模型，而是把一次全屏读取改成多次局部搜索：先由规划模型推测目标可能所在的区域，再裁剪放大交给定位模型；它改进的是模型看到的观测。

\Needspace{60mm}
\subsubsection{SeeClick}
\label{sec:vision-understanding-seeclick}

此前的界面智能体多依赖网页HTML或系统无障碍树等结构化文本来确定控件，这些信息在许多应用中不可得，或者冗长难用；而通用视觉语言模型虽然能读截图，却很少学过把指令落到精确坐标。SeeClick的主要想法是只用截图，并把坐标写成文字，让语言模型像生成普通答案一样生成位置；所需的大量训练记录则从网页和移动界面自动收集\scite{vision-seeclick2024}。

模型以Qwen-VL为基础，结构沿用不变，各组件规模见表\ref{tab:vision-understanding-seeclick-config}\scite{vision-qwenvl2023}。截图缩放到$448\times448$，视觉编码器是以OpenCLIP的ViT-bigG初始化的ViT，按步长14切块，得到$32\times32=1024$个图像块特征；适配器是一层随机初始化的交叉注意力，用256个可学习查询读取这些特征，并在查询与键上加入二维绝对位置编码，把视觉信息压缩为256个单元；这些单元放在指令文字前面送入以Qwen-7B初始化的语言模型。输出是一段文字，例如\code{click (0.49, 0.40)}，两个数是保留两位小数、取值在$[0,1]$的比例，分别表示相对于截图宽、高的位置。这与Qwen-VL本身用$[0,1000)$整数表示框坐标的约定不同；SeeClick把数字当作普通文字处理，不增加专门的位置词元。由于位置以文字形式生成，模型不需要新增回归头，指令跟随和坐标预测共用同一个解码器。

\begin{table}[htbp]
  \centering
  \small
  \begin{tabularx}{\linewidth}{@{}p{16mm}XXp{12mm}@{}}
    \toprule
    组件 & 结构 & 输入与输出 & 参数量\\
    \midrule
    视觉编码器 & ViT，以OpenCLIP ViT-bigG初始化，切块步长14 & $448\times448$截图 \ensuremath{\to} 1024个图像块特征 & 1.9B\\
    适配器 & 单层交叉注意力，256个可学习查询，二维绝对位置编码 & 1024个特征 \ensuremath{\to} 256个视觉单元 & 0.08B\\
    语言模型 & 以Qwen-7B初始化的解码器 & 256个视觉单元与指令词元 \ensuremath{\to} 坐标文字 & 7.7B\\
    合计 & — & — & 9.6B\\
    \bottomrule
  \end{tabularx}
  \caption{SeeClick所用Qwen-VL的组件规模。参数量取自Qwen-VL论文，合计为原文给出的约数；视觉编码器的层数与宽度沿用ViT-bigG，SeeClick论文未另行给出。}
  \label{tab:vision-understanding-seeclick-config}
\end{table}

训练数据约100万条，混合三类来源：从约30万个网页自动提取的可见文字及其悬停提示与对应位置，约43万条；移动界面的控件描述、界面摘要等数据，约42万条；以及约14.5万条LLaVA一般视觉语言指令数据，用于保持原有能力。每条记录可以组织成“文字找位置”或“位置说文字”两种方向。监督是目标输出序列的交叉熵，模型逐词预测位置文本，不是用欧氏距离回归坐标。训练时用LoRA适配视觉编码器和语言模型，即在原权重旁加入低秩增量矩阵，只更新这些增量；约训练1万步，相当于一轮，全局批大小64，初始学习率$3\times10^{-5}$并按余弦衰减。论文同时建立了ScreenSpot评测集，含600余张截图、1200余条指令，覆盖iOS、Android、macOS、Windows和网页界面。

使用时，把截图和指令按训练格式输入，解析输出中的两个数，检查是否落在$[0,1]$之内，再乘以截图尺寸。若截图宽1600、高1000像素，输出$(0.49,0.40)$对应连续坐标$(784,400)$；实际点击还需与显示缩放和截图坐标一致。数字生成损失并不直接等于像素误差，“0.49”与“0.94”只差一个词元的顺序，像素距离却很远。定位评价可检查点是否落在目标控件框内；点落在正确按钮上，不能说明该按钮当前可用，更不能说明保存已经成功，多步操作还需在新截图中确认结果。

\subsubsection{ScreenSeekeR}
\label{sec:vision-understanding-screenseeker}

SeeClick一次读取整张截图，输入分辨率受模型限制。专业软件的高分辨率屏幕上，目标控件只占屏幕面积的极小比例，压缩到固定输入尺寸时可能只剩几个像素。ScreenSpot-Pro为此收集了专业应用的高分辨率截图作为评测基准：1581条指令，每条对应一张不同的截图，来自23个应用，分辨率多在1080p以上，目标平均只占截图面积的0.07\%；同一论文提出的ScreenSeekeR是一种搜索方法，二者名称不能互换\scite{vision-screenseeker2025}。ScreenSeekeR的主要想法是利用界面的层级布局知识缩小搜索区域：保存按钮通常在工具栏，工具栏通常在窗口上方，于是可以先找较大的区域，再逐级放大。

方法组合两个已有模型，不做训练。原实现以GPT-4o作为规划模型，以OS-Atlas-7B作为定位模型。规划模型读取截图和指令，用文字给出目标本身以及可能邻近的控件或区域；定位模型把这些文字描述落到框。候选区域的评分不直接采用规划模型的口头置信度，而是让定位框的中心为候选区域投票：将中心$(x,y)$转换为区域内部归一化位置$(x',y')$，在区域内时给予
\begin{equation}
  s=\exp\!\left(-\frac{(x'-0.5)^2+(y'-0.5)^2}{2\sigma^2}\right)
  \label{eq:vision-screenseeker-vote}
\end{equation}
的贡献，区域外贡献为0，原设置取$\sigma=0.3$。将所有定位框的贡献相加，再用非极大值抑制减少重复候选。靠近区域中心的多条线索可以提高该区域的优先级，但这种几何投票仍可能支持错误窗口。

搜索按评分依次进入裁剪区域，在其中继续提出位置或进行直接定位，直到规划模型确认找到目标、区域小于所定尺寸或到达最大深度。原设置在区域足够小时直接调用定位模型，这一尺寸超参数取1280像素；最小区域尺寸与最大搜索深度作为算法输入出现，论文未给出具体数值。这种逐级缩小的搜索，正是\ref{sec:vision-understanding-challenge-resolution}节所说多尺度观察的延续：经典方法在图像金字塔的每一层上滑窗，这里改由规划模型根据布局知识决定只放大哪几个区域。每次裁剪都应保存到父图的坐标变换，再组合回全屏。例如，父裁剪左上角为$(400,200)$，其内再裁出左上角$(100,50)$、宽高$200\times100$的区域；若最终归一化点为$(0.25,0.60)$，全屏点为$(400+100+0.25\times200,\ 200+50+0.60\times100)=(550,310)$。若中途调整过尺寸，还要先逆转该次缩放。

在ScreenSpot-Pro上，同一个OS-Atlas-7B定位模型单次读取全屏时的平均准确率为18.9\%，接入ScreenSeekeR的搜索后为48.1\%。两个数字使用同一定位模型和同一评测集，差别来自观测方式的改变，这与“小控件在整屏缩放中丢失”这一失败解释一致。

两种做法的比较见表\ref{tab:vision-understanding-gui-compare}。ScreenSeekeR保留了定位模型的单次读取方式，把预算从一次全屏读取改为多次局部搜索，换来了小控件的有效分辨率，代价是多次调用大模型的时间与费用，以及对规划模型布局知识的依赖。裁剪虽然放大了目标，也可能删去“当前窗口”等消歧线索，因此递归状态仍需保留整屏上下文与候选所属位置。两者都只对当前截图产生定位，历史动作和界面状态变化应作为下一次观测重新进入系统。

\begin{table}[htbp]
  \centering
  \small
  \begin{tabularx}{\linewidth}{@{}p{16mm}XX@{}}
    \toprule
    比较项 & SeeClick & ScreenSeekeR\\
    \midrule
    改进对象 & 模型的定位能力 & 模型看到的观测\\
    训练 & 界面定位数据上LoRA适配 & 不训练，组合已有模型\\
    每条指令的调用 & 一次生成 & 多次规划、定位与裁剪\\
    高分辨率处理 & 整屏缩放到固定输入 & 逐级裁剪放大\\
    主要失败 & 小控件在缩放中丢失 & 规划误导、裁剪丢失上下文\\
    \bottomrule
  \end{tabularx}
  \caption{两种界面定位做法的对照。两者输出都是截图坐标，评价口径都是点是否落在目标控件内。}
  \label{tab:vision-understanding-gui-compare}
\end{table}

\section{文字读取与结构化解析}
\label{sec:vision-understanding-reading}

定位找到了票据区域，但“看见数字”和“知道数字属于哪个字段”是两回事。表格中，即使所有字符都读对，合并表头归属错误仍会改变含义；图表中，柱子的高度只有对照刻度和图例才成为数值。本节的输入是含文字或符号的图像，输出是字符串以及它们之间的关系：阅读顺序、表格的行列、字段的键值、图表的数据记录、图解的连接。

本节按被恢复的对象展开，对象的结构由简到繁。字符串的读取是其余各项的基础，要回答“这里写了什么”。页面结构的解析在字符串之上恢复区域类型和阅读关系。表格、表单、图表和图解各有专门的结构：表格靠行列与跨格，表单靠键与值，图表靠几何编码与刻度，图解靠节点与有向连线。后四种对象的恢复大多复用前两项的模型，再加上对象特有的结构规则，因此不再按模型细分。

\subsection{字符串的读取}
\label{sec:vision-understanding-ocr}

\term{光学字符识别}（optical character recognition, OCR）把文字图像转换为机器可处理的字符串。输入是一幅场景图或页面图，输出是若干文字实例，每个实例带一个框或多边形以及一个字符串。完整流程通常分为两步：文字检测找出哪些区域包含文字，文字识别把每个区域读成字符串。位置正确而内容读错、内容正确但属于另一行，是两种不同错误，分别出现在这两步。

经典的文档OCR用\ref{chap:vision-image-features}章的工具完成这两步。检测阶段先二值化：文档中墨迹与纸张的灰度分布通常呈两个峰，Otsu方法（\ref{sec:vision-features-otsu}节）选择使两类类间方差最大的阈值，光照不均时改用局部自适应阈值；随后做连通分量标记（\ref{sec:vision-formation-connectivity}节），每个分量大致是一个字符或字符的一部分，再按分量的位置和大小把它们聚成行，例如沿水平方向统计墨迹像素数，用投影的谷值切分行与字。识别阶段把每个字符图像缩放到固定大小，与字库模板做归一化互相关（\ref{sec:vision-features-template-matching}节），或提取HOG等梯度特征（\ref{sec:vision-features-hog}节）后交给分类器。这条流水线在白纸黑字、版式规整的扫描件上有效，因为它依赖两个条件：前景与背景可以用一个阈值分开，字符彼此分离而且排成直线。场景文字两个条件都不满足：招牌背景复杂，阈值分不开文字与纹理；字符可能粘连、弯曲或带透视变形，连通分量与直线投影都会失效。下面的两种学习模型保留了这条流水线“先找部件、再组合、再识别”的结构，但把每个部件换成了从数据中学到的判断。

两步各有难点。检测要处理弯曲、倾斜、长短悬殊的文字行，直接回归一个矩形框可能包含大量背景，也可能把一行切断。识别要处理模糊、手写和罕见字符串，还要决定语言上下文能在多大程度上帮助判断字形。CRAFT回答检测的难点：先找单个字符，再判断相邻字符是否相连，由下而上拼出任意形状的文字。TrOCR回答识别的难点：用预训练的图像Transformer读取文字行，再用预训练的文字Transformer逐词生成字符串。两者处在流程的不同环节，常被串联使用。

\subsubsection{CRAFT}
\label{sec:vision-understanding-craft}

弯曲招牌中的字符未必落在同一条直线上，以整词为单位回归框的检测器难以给出贴合的边界；长文字行的感受野也可能不够覆盖整行。CRAFT的主要想法是把检测拆成两个局部判断：哪里是字符中心，哪里是两个相邻字符之间的连接。局部判断只需看字符尺度的邻域，组合起来却能描出任意长度、任意弯曲的文字行\scite{vision-craft2019}。

模型是一个全卷积网络。输入为任意尺寸的$H\times W\times3$图像，编码器采用带批归一化的VGG-16，解码器逐级上采样并通过跳跃连接融合编码器各层的特征，结构类似U-Net；卷积网络的基本运算可参考\BookChapterRef{chap:convolutional-networks}{卷五的卷积神经网络章}。输出是两张稠密图：\emph{字符区域图}$S_{\mathrm r}$给出每个位置作为字符中心的程度，\emph{关联图}$S_{\mathrm a}$给出每个位置作为相邻字符连接处的程度。

论文只给出了骨干和跳跃连接的示意，表\ref{tab:vision-understanding-craft-layers}的逐层配置依据作者发布的官方实现。编码器取VGG-16-BN的前若干卷积组，在四个分辨率上取出跳跃特征$\bm s_1,\ldots,\bm s_4$；末端用一个步长为1的最大池化和一个扩张率为6的$3\times3$卷积代替原VGG的全连接层fc6，再接$1\times1$卷积代替fc7，在不降低分辨率的情况下扩大感受野，使长文字行的上下文也能被看到。解码器有四个上采样块，每块的计算顺序是：把上一级输出用双线性插值放大到下一个跳跃特征的尺寸，与该跳跃特征沿通道拼接，经$1\times1$卷积、批归一化、ReLU压缩通道，再经$3\times3$卷积、批归一化、ReLU输出。第一块的输入不需要放大，直接拼接fc7的输出与$\bm s_4$。最后的分类头由三层$3\times3$卷积和两层$1\times1$卷积组成，输出两个通道，分别作为$S_{\mathrm r}$与$S_{\mathrm a}$。由于最浅的跳跃特征位于$1/2$分辨率，输出图的边长是输入的一半。

\begin{table*}[tbp]
  \centering
  \small
  \begin{tabularx}{\linewidth}{@{}p{24mm}Xp{24mm}p{14mm}p{10mm}p{22mm}@{}}
    \toprule
    阶段或层组 & 层类型 & 卷积核与步长 & 输出通道 & 重复 & 输出尺寸\\
    \midrule
    卷积组1 & 卷积、批归一化、ReLU & $3\times3$，1 & 64 & 2 & $H\times W$\\
    卷积组2 & 最大池化；卷积、批归一化、ReLU & $2\times2$，2；$3\times3$，1 & 128 & 1；2 & $H/2$，取出$\bm s_1$\\
    卷积组3 & 最大池化；卷积、批归一化、ReLU & $2\times2$，2；$3\times3$，1 & 256 & 1；2 & $H/4$，取出$\bm s_2$\\
    卷积组4 & 卷积；最大池化；卷积、批归一化、ReLU & $3\times3$，1；$2\times2$，2；$3\times3$，1 & 256；512 & 1；1；2 & $H/8$，取出$\bm s_3$\\
    卷积组5 & 卷积；最大池化；卷积、批归一化、ReLU & $3\times3$，1；$2\times2$，2；$3\times3$，1 & 512；512 & 1；1；2 & $H/16$，取出$\bm s_4$\\
    代替fc6、fc7 & 最大池化；扩张卷积；卷积 & $3\times3$，1；$3\times3$扩张6，1；$1\times1$，1 & 1024 & 各1 & $H/16$\\
    上采样块1 & 拼接fc7输出与$\bm s_4$；卷积；卷积 & $1\times1$；$3\times3$ & 512；256 & 1 & $H/16$\\
    上采样块2 & 放大并拼接$\bm s_3$；卷积；卷积 & $1\times1$；$3\times3$ & 256；128 & 1 & $H/8$\\
    上采样块3 & 放大并拼接$\bm s_2$；卷积；卷积 & $1\times1$；$3\times3$ & 128；64 & 1 & $H/4$\\
    上采样块4 & 放大并拼接$\bm s_1$；卷积；卷积 & $1\times1$；$3\times3$ & 64；32 & 1 & $H/2$\\
    分类头 & 卷积与ReLU；卷积与ReLU；卷积 & $3\times3$；$1\times1$；$1\times1$ & 32、32、16；16；2 & 3；1；1 & $H/2\times W/2\times2$\\
    \bottomrule
  \end{tabularx}
  \caption{CRAFT的逐层结构，依据官方实现。输出尺寸只写高度，宽度按相同比例变化；例如输入$768\times768$时，输出为$384\times384\times2$。官方实现在原VGG第2组最后一个卷积之后、第3至5组倒数第二个卷积之后取出跳跃特征（取在批归一化之后、ReLU之前）；组内剩余的卷积在下一段开头执行，表中按执行顺序写在下一行，原VGG第5组的最后一个卷积未被使用。}
  \label{tab:vision-understanding-craft-layers}
\end{table*}

训练目标需要两张图的标注，而大多数真实数据只有词级框。合成数据集SynthText提供字符框，可把二维高斯模板通过透视变换映射到每个字符框内，形成区域目标$S_{\mathrm r}^*$；相邻两个字符框各取上下三角形中心，连成的四边形内再放高斯模板，形成关联目标$S_{\mathrm a}^*$。对只有词级框和字符串的真实数据，原方法用当前模型自己生成字符级伪标注：裁出词区域，预测区域图，用分水岭算法（\ref{sec:vision-features-watershed}节）切分出若干字符框，再把坐标还原到原图。分水岭把区域图看成地形，每个字符中心是一个盆地，从盆地向外“注水”直到相邻盆地相遇，相遇处就是字符之间的分界。

伪标注可能切错，因此要衡量可靠性。设一个词的真实长度为$l(w)$，切分得到的字符数为$l^{\mathrm c}(w)$，置信度定义为
\begin{equation}
  s_{\mathrm{conf}}(w)=\frac{l(w)-\min\bigl(l(w),\,|l(w)-l^{\mathrm c}(w)|\bigr)}{l(w)}\text{。}
  \label{eq:vision-craft-conf}
\end{equation}
切出的字符数与真实长度一致时置信度为1，相差越多越低。若$s_{\mathrm{conf}}<0.5$，就放弃切分结果，改为按字符数均分词区域，并把置信度置为0.5。词内各像素取所属词的置信度，组成置信图$S_{\mathrm c}$，合成数据上取1。训练损失为逐像素加权的平方误差
\begin{equation}
  L=\sum_{p}S_{\mathrm c}(p)\Bigl(\bigl\lVert S_{\mathrm r}(p)-S_{\mathrm r}^*(p)\bigr\rVert_2^2
   +\bigl\lVert S_{\mathrm a}(p)-S_{\mathrm a}^*(p)\bigr\rVert_2^2\Bigr)\text{。}
  \label{eq:vision-craft-loss}
\end{equation}
置信度低的伪标注贡献较小，模型不会被明显切错的字符框带偏。原论文先在SynthText上训练5万次迭代，再在各真实数据集上微调2.5万次迭代，微调时按1比5的比例混入合成数据，并以1比3的正负像素比例做困难样本挖掘；各阶段都用Adam优化。论文没有给出学习率、批大小和训练图像尺寸。

推理时，将区域图超过阈值$\tau_{\mathrm r}$或关联图超过阈值$\tau_{\mathrm a}$的位置标为1，得到二值图；做连通分量标记（\ref{sec:vision-formation-connectivity}节），每个分量就是一个文字实例，再为其拟合最小面积的旋转矩形，或沿字符中心构造适合弯曲文字的多边形。论文没有给出两个阈值的取值，官方实现默认以0.4作为区域图和关联图的二值化阈值，并只保留区域图最大值超过0.7的分量。这一步与经典流水线中对二值化墨迹做连通分量分组的做法相同，区别在于被分组的是学到的字符响应和关联响应：字符之间即使在像素上粘连或断开，只要关联图的响应正确，分组就仍然正确。教学招牌“青山路”若三个字符响应之间有两条正确关联，会形成一个文字实例；若“路”与旁边门牌之间也出现强关联，分组就可能错误合并。此时问题发生在文字检测，后续识别器再准确也无法恢复正确的实例边界。

\subsubsection{TrOCR}
\label{sec:vision-understanding-trocr}

文字区域确定后，还要把裁出的文字图像读成字符串。此前的识别器通常以卷积网络提取特征，接循环网络逐列建模，再用连接时序分类等方法输出字符，需要时另接外部语言模型。这种结构的视觉与语言部分难以直接利用大规模预训练的Transformer。TrOCR的主要想法是把识别写成图像到文字的序列转换：编码器和解码器分别用已有的图像Transformer和文字Transformer初始化，二者之间只新增交叉注意力，由识别数据学会连接\scite{vision-trocr2022}。

输入文字图像$\bm X_{\mathrm{text}}$被缩放到$384\times384$，切成$16\times16$的图像块，共$24\times24=576$块；每块展平后线性投影为$d$维，加上位置表示，送入从BEiT或DeiT初始化的ViT编码器，输出$576\times d$的视觉表示。解码器从RoBERTa或MiniLM初始化，按标准编码器—解码器结构读取已生成的词元，并以交叉注意力读取视觉表示，预测下一个子词词元；交叉注意力的计算可参考\BookChapterRef{chap:transformer}{卷五的Transformer章}。这一固定尺寸设置把长文字行压成正方形，与后文Pix2Struct保持纵横比的设计不同。

论文给出了三种规模，见表\ref{tab:vision-understanding-trocr-config}。三种规模的输入都是$384\times384$、图像块$16\times16$，编码器序列为576个图像块，另加一个[CLS]词元，用DeiT初始化时再加一个蒸馏词元。解码器层数少于用来初始化的文字模型时，取该模型的后一半层，例如从24层的RoBERTa-Large取后12层；文字模型原本只有编码器，没有交叉注意力，这部分参数随机初始化。

\begin{table*}[tbp]
  \centering
  \small
  \begin{tabularx}{\linewidth}{@{}p{16mm}p{14mm}p{12mm}Xp{18mm}Xp{14mm}@{}}
    \toprule
    规模 & 图像块 & 词元数 & 编码器（初始化；层数、宽度、头数） & 解码器初始化 & 解码器（层数、宽度、头数） & 参数量\\
    \midrule
    TrOCR-Small & $16\times16$ & 576 & DeiT-Small；12层，384，6头 & MiniLM & 6层，256，8头 & 62M\\
    TrOCR-Base & $16\times16$ & 576 & BEiT-Base；12层，768，12头 & RoBERTa-Large & 12层，1024，16头 & 334M\\
    TrOCR-Large & $16\times16$ & 576 & BEiT-Large；24层，1024，16头 & RoBERTa-Large & 12层，1024，16头 & 558M\\
    \bottomrule
  \end{tabularx}
  \caption{TrOCR的三种规模，依据原论文的实验设置。词元数为编码器的图像块数，不含特殊词元；MLP倍率论文未单独给出，沿用各初始化模型。}
  \label{tab:vision-understanding-trocr-config}
\end{table*}

编码器的单个块采用ViT的写法（\ref{sec:vision-ivm-vit}节），层归一化放在子层之前：
\[
  \bm z'=\bm z+\operatorname{MSA}(\operatorname{LN}(\bm z)),\qquad
  \bm z''=\bm z'+\operatorname{MLP}(\operatorname{LN}(\bm z'))\text{。}
\]
解码器的单个块依次做三步：带因果掩码的自注意力，使第$j$个位置只能看到已生成的词元；以编码器输出为键和值的交叉注意力；前馈网络。每步都带残差连接和层归一化。文字图像的信息只通过交叉注意力进入解码器，而解码器的自注意力承载语言上下文，这两个通道的相对强弱决定了后文所说语言先验的影响程度。

以文字图像$\bm X_{\mathrm{text}}$和目标词元$y_1,\ldots,y_m$为一条记录，训练最小化
\begin{equation}
  L_{\mathrm{OCR}}
  =-\sum_{j=1}^{m}\log p_{\bm\theta}
    (y_j\mid y_{<j},\bm X_{\mathrm{text}})\text{。}
  \label{eq:vision-trocr-loss}
\end{equation}
训练分阶段进行：第一阶段在合成的印刷文字行上预训练，共6.84亿行；第二阶段用规模较小、分别面向任务构建的数据继续预训练，印刷体约330万行，手写体约1790万行，场景文字则用MJSynth与SynthText约1600万幅；最后在票据（SROIE）、手写（IAM）和场景文字基准上微调。全部参数参与更新，批大小2048，学习率$5\times10^{-5}$。输出用BPE与SentencePiece分词，得到子词词元。

使用时从起始符逐词解码，直到结束符或长度上限；用束搜索保留若干候选，原设置的束宽为10，最终去掉特殊词元并恢复字符串。这里的预测单元是子词，一步未必恰好对应一个字符。模糊字形“1”和“I”可能在视觉上相近，语言上下文会影响选择：对常见地名，这有时有助于消歧；对票据编号或未见专名，强语言先验也可能把正确的罕见字符串改成常见词。因此应保留原裁剪和识别结果用于核对。

CRAFT与TrOCR分工明确，比较时要放在各自的环节中看。CRAFT以字符和连接为单位，换来了对弯曲和长文字行的适应，代价是依赖伪标注质量，且字符间距很大或很小时容易断开或误连；它不输出字符串。TrOCR把识别交给预训练的编码器和解码器，换来了语言上下文和手写体的处理能力，代价是固定输入尺寸和语言先验可能改写罕见字符串；它不给出全页文字位置、双栏阅读顺序或表格归属。二者串联时，检测阶段的合并或截断错误会原样传给识别阶段。

\subsection{页面结构的解析}
\label{sec:vision-understanding-page}

一页文档上的文字可能属于标题、正文、表格、脚注或公式。把左栏第一行和右栏第一行交错拼接，即使每个字都正确，仍会得到错误句子。这一小节的输入是页面图像，输出是带类型的页面元素、文字所属的字段或区域，以及元素之间的阅读关系。难点在于：语义既依赖文字内容，也依赖二维位置和视觉样式，三者需要在同一个表示中联系起来。

两条路线对“文字从哪里来”给出不同回答。LayoutLMv3依赖外部OCR提供词及其框，把文字、二维布局和图像块共同编码，再接分类或抽取任务头输出结构；它擅长在已读出的文字上做判断。Pix2Struct不使用外部OCR，直接从像素生成结构化文字序列；它省去OCR环节，但要求视觉编码器自己保留小字。

\subsubsection{LayoutLMv3}
\label{sec:vision-understanding-layoutlm}

此前的文档预训练模型通常用卷积网络或检测器提取页面图像特征，并为文字和图像设计不同的预训练目标，例如文字一侧做掩码语言建模，图像一侧做像素重建。两种目标的粒度不一致，难以让模型学到“这个词对应页面上哪块图像”。LayoutLMv3的主要想法是统一掩码：文字和图像块都用离散单元表示、都做掩码预测，再加一个专门的词块对齐目标，把文字与它所在的图像块联系起来；图像一侧也不再使用卷积骨干，只做线性投影\scite{vision-layoutlmv32022}。

文字一侧的输入来自外部OCR：每个词经分词得到词元，词元嵌入由RoBERTa初始化，再加上一维序列位置和二维布局位置。布局位置把所属文字段的框坐标按页面尺寸归一化，分别嵌入横纵坐标与宽高；同一文字段内的词共享一个框。图像一侧把页面缩放到$224\times224$，切成$16\times16$的块，共196块，各块线性投影后加上一维位置嵌入。两组单元拼接后送入Transformer，输出每个词元和每个图像块的上下文表示。这样，“128.50”的表示可以同时参考附近的“总额”及其页面位置。表\ref{tab:vision-understanding-layoutlmv3-config}给出两种规模的配置。文字序列至多512个词元，用BPE分词；嵌入层和Transformer以RoBERTa的权重初始化。图像块只加一维可学习位置嵌入，论文报告加二维位置嵌入没有带来改善。

\begin{table*}[tbp]
  \centering
  \small
  \begin{tabularx}{\linewidth}{@{}p{22mm}p{12mm}Xp{8mm}p{10mm}p{8mm}p{14mm}p{12mm}@{}}
    \toprule
    规模 & 图像块 & 词元数 & 层数 & 宽度 & 头数 & MLP倍率 & 参数量\\
    \midrule
    LayoutLMv3-Base & $16\times16$ & 至多512个文字词元，196个图像块 & 12 & 768 & 12 & 4（3072） & 133M\\
    LayoutLMv3-Large & $16\times16$ & 至多512个文字词元，196个图像块 & 24 & 1024 & 16 & 4（4096） & 368M\\
    \bottomrule
  \end{tabularx}
  \caption{LayoutLMv3的两种规模，依据原论文。图像输入为$224\times224$；MLP倍率一栏括号内为前馈层中间维数。}
  \label{tab:vision-understanding-layoutlmv3-config}
\end{table*}

单个Transformer块沿用RoBERTa的结构，在残差相加之后做层归一化，但自注意力的打分加入了位置偏置。设第$i$、$j$个单元的查询与键为$\bm q_i$、$\bm k_j$，每个头的维数为$d_h$，打分为
\[
  e_{i,j}=\frac{\bm q_i^\top\bm k_j}{\sqrt{d_h}}+b^{\mathrm{1D}}_{i,j}+b^{\mathrm{2D}}_{i,j}\text{，}
\]
其中$b^{\mathrm{1D}}_{i,j}$由两个单元在序列中的相对距离决定，$b^{\mathrm{2D}}_{i,j}$由两个文字段在页面上的横、纵相对距离决定，二者都是可学习的偏置，沿用LayoutLMv2的做法。随后$\bm Z'=\operatorname{LN}(\bm Z+\operatorname{MSA}(\bm Z))$，$\bm Z''=\operatorname{LN}(\bm Z'+\operatorname{FFN}(\bm Z'))$。二维相对偏置让“同一行右侧相邻”这类版面关系直接影响注意力，而不必完全依靠绝对坐标嵌入去推断。

预训练使用约1100万页IIT-CDIP文档图像，目标由三项相加：
\begin{equation}
  L=L_{\mathrm{MLM}}+L_{\mathrm{MIM}}+L_{\mathrm{WPA}}\text{。}
  \label{eq:vision-layoutlmv3-loss}
\end{equation}
文字掩码项$L_{\mathrm{MLM}}$以片段方式遮去30\%的文字词元，片段长度服从参数为3的泊松分布，用交叉熵预测被遮词元。图像掩码项$L_{\mathrm{MIM}}$按块状区域遮去约40\%的图像块，预测这些块的离散视觉编码；编码由一个离散变分自编码器给出，它以文档图像模型DiT的图像分词器初始化，词表含8192个视觉词元，而不是直接回归原始像素。\term{词块对齐}（word-patch alignment, WPA）项只作用于未被遮的文字词元：若该词对应的图像块也未被遮，标为“对齐”，否则标为“未对齐”，由一个两层感知机以二元交叉熵预测。预训练共50万步，批大小2048。若一行金额文字位于图像块7和8，遮去这两块而保留词元时，对齐标签就变为“未对齐”；模型要答对，必须知道这行文字在页面上的位置对应哪些图像块。这个目标利用OCR位置建立的联系，不是判断任意一段文字与整幅图是否语义一致，也不能保证OCR本身读对。

下游使用时在预训练表示上加任务头并微调全部参数。表单理解按词元或文字段预测“问题”“答案”“标题”“其他”等类别；票据信息抽取预测字段标签；文档问答在词元上预测答案的起止位置；文档分类读取整体表示。版面分析是另一种配置：不输入文字，只把LayoutLMv3的图像部分作为检测骨干，从基础版的第4、6、8、12层取出特征，经改变分辨率的模块变成多尺度特征，再接带特征金字塔的Cascade R-CNN检测器，输出标题、表格、图片等区域框。两种配置的输入不同，LayoutLMv3因此不能被理解为一个无条件输出整页HTML及阅读顺序的解码器；能恢复什么，由输入、任务头和监督决定。

\subsubsection{Pix2Struct}
\label{sec:vision-understanding-pix2struct}

LayoutLMv3式方法在使用时必须先运行OCR，OCR的错误和耗时都会传入下游；同时，把页面缩放到固定正方形会让长条网页或宽表格的小字失真。Pix2Struct的主要想法有两点：一是不用OCR，直接从像素生成结构化文字；二是按输入预算可变地缩放图像，尽量保留原纵横比。它的预训练任务是把网页截图解析成简化的HTML，网页自带的结构为“从像素到结构”提供了大量免费监督\scite{vision-pix2struct2023}。

编码器是ViT。给定最多$n_{\max}$个$16\times16$图像块的预算，图像按原纵横比放大或缩小，使切出的块数不超过预算且尽可能多；块的位置用行号和列号两套二维绝对位置嵌入表示，因此不同形状的页面都能被正确编码。原论文预训练时预算为2048块。解码器是Transformer文字解码器，读取编码器输出并逐词生成目标序列。论文把这种结构称为基于ViT的图像编码器—文字解码器，并把整体类比为视觉上的T5：各种任务都化为“输入图像、输出文字”。两种规模的配置见表\ref{tab:vision-understanding-pix2struct-config}。

\begin{table*}[tbp]
  \centering
  \small
  \begin{tabularx}{\linewidth}{@{}p{22mm}p{12mm}Xp{16mm}p{10mm}p{10mm}p{12mm}p{12mm}@{}}
    \toprule
    规模 & 图像块 & 编码器块数上限 & 层数 & 宽度 & 头数 & MLP倍率 & 参数量\\
    \midrule
    Pix2Struct-Base & $16\times16$ & 预训练2048；微调4096（信息图问答6144） & 编码器12，解码器12 & 768 & 论文未给出 & 论文未给出 & 282M\\
    Pix2Struct-Large & $16\times16$ & 预训练2048；微调3072 & 18（编码器与解码器的划分论文未给出） & 1536 & 论文未给出 & 论文未给出 & 1.3B\\
    \bottomrule
  \end{tabularx}
  \caption{Pix2Struct的两种规模，依据原论文。块数上限是可变分辨率输入的预算，实际块数随页面纵横比变化；预训练时解码器序列长度为128个词元。}
  \label{tab:vision-understanding-pix2struct-config}
\end{table*}

可变分辨率的计算可以用一个教学例说明。设页面为$1600\times900$像素，预算为2048块。保持纵横比缩放$\alpha$倍后，块网格为$\lfloor1600\alpha/16\rfloor\times\lfloor900\alpha/16\rfloor$；要求乘积不超过2048，约有$100\alpha\times56.25\alpha\leq2048$，即$\alpha\leq0.603$。取$\alpha=0.6$得到$60\times33=1980$块，页面被缩到约$960\times540$像素。若改用$224\times224$的固定正方形，横向缩小约7倍、纵向缩小约4倍，字形还被拉扁。二维位置嵌入用行号和列号分别查表，所以$60\times33$与$33\times60$两种网格都能被正确编码。

预训练数据来自C4语料中网址对应的约8000万张网页截图及其HTML。原方法保留可见内容相关的节点，压缩冗长的嵌套，只保留节点的文字与图像替代信息，再选取能放入目标长度的子树；截图上标出该子树的区域，并以随机片段遮去其中一半文字，训练解码器生成完整的简化结构。训练目标是目标序列的逐词交叉熵，形式与式\eqref{eq:vision-trocr-loss}相同。直接从解析任务开始训练不稳定，因此先做阅读热身：把BooksCorpus的文本以随机字体和颜色渲染成图，输入上限128块，训练3万步读出文字，再进入截图解析。截图解析阶段，基础版再训练27万步、批大小2048，大型版训练17万步、批大小1024。

下游任务以同一接口微调。问题等文字条件不作为单独的文字输入，而是渲染在图像顶部作为标题区，与页面一同被编码；指代任务则把框直接画在截图上。论文在四个领域的九个任务上微调：插图类的ChartQA、AI2D、OCR-VQA，界面类的RefExp、Widget Captioning、Screen2Words，自然图像类的TextCaps，以及文档类的DocVQA和InfographicVQA。使用时必须遵循相应检查点的预处理和输出约定。若教学页面具有“发票—编号A017—总额128.50”的层次，可以用嵌套节点表示从属关系；解析器应检查节点是否闭合、文字是否遗漏及金额属于哪个字段。字符串符合语法，只说明可以读入结构，不说明内容忠实于页面。改变块预算会改变可保留的视觉细节，解码更长的字符串不能恢复已经因缩小而不可辨认的小数点。

页面元素恢复后，还要按既定关系组织阅读顺序。双栏正文应在本栏顺读，再转入下一栏；跨栏标题和所属公式要与相应内容连接。OmniDocBench同时提供完整文档解析和文字、版面、表格、公式等模块的评价\scite{vision-omnidocbench2025}。这些分项使整页误差能够定位到具体结构，而非仅用一串文字相似度概括。

两种路线的对照见表\ref{tab:vision-understanding-page-compare}。Pix2Struct把LayoutLMv3依赖的外部OCR换成了模型自己的读取，把固定正方形输入换成了可变分辨率，换来了无需OCR的单一接口和对网页结构的利用；代价是所有文字都要从像素中读出，小字能否保留完全取决于块预算，且生成的结构串需要额外的语法与内容检查。LayoutLMv3在OCR可靠时能直接利用准确的文字和坐标，输出也天然对应到具体词元，便于追溯来源。

\begin{table}[htbp]
  \centering
  \small
  \begin{tabularx}{\linewidth}{@{}p{14mm}XX@{}}
    \toprule
    比较项 & LayoutLMv3 & Pix2Struct\\
    \midrule
    文字来源 & 外部OCR的词与框 & 模型从像素读取\\
    图像输入 & 固定$224\times224$ & 按块预算可变缩放，保持纵横比\\
    预训练 & 文字与图像块掩码预测、词块对齐 & 截图到简化HTML的生成\\
    输出方式 & 任务头：分类、跨度、检测 & 自回归生成结构化文字\\
    来源追溯 & 输出对应到OCR词元及其框 & 需另行对齐生成内容与页面\\
    主要失败 & OCR错误直接传入 & 缩放丢失小字、结构串不忠实\\
    \bottomrule
  \end{tabularx}
  \caption{两种页面解析路线的对照。比较以同一页面、同一结构化任务为前提。}
  \label{tab:vision-understanding-page-compare}
\end{table}

\subsection{表格结构的解析}
\label{sec:vision-understanding-tables}
\label{sec:vision-understanding-pubtables}

表格中，一个数的含义来自它所在的行、列和表头。输入是页面或裁剪的表格图像，输出是逻辑表：每个单元的行列索引、跨度、内容和来源框。难点在于行列边界往往没有画线，表头可能跨多行多列，仅按OCR字符串的出现顺序排列会丢失这些关系。

有框线的表格可以用经典方法恢复。二值化后，用一个很长的水平线段作结构元素做形态学开运算（\ref{sec:vision-transforms-binary-morphology}节），只有长度超过结构元素的水平笔画能够保留，文字笔画被去掉，剩下的就是横线；竖线同理。也可以在边缘图上用Hough变换（\ref{sec:vision-features-hough}节）检测直线：每个边缘点在参数空间中为所有经过它的直线投票，票数集中的参数就是图中的长直线。横线与竖线的交点给出单元格的角点。这类方法在框线完整时可靠，但许多表格只有表头下的一条横线，甚至完全没有框线，行列边界只能从文字的对齐和间距推断；跨格表头也不会在框线上留下额外线索。学习方法改为直接预测行、列和跨格的区域，不再依赖框线是否存在。

PubTables-1M的做法是把表格结构当作一组可检测的对象\scite{vision-pubtables2022}。此前的表格数据集标注常不一致，例如同一个跨两列的表头，有的标为一格，有的标为一格文字加一格空白；模型因而学到不一致的结构。该工作从PubMed Central开放获取论文中收集近百万张表格，并对标注做\emph{规范化}：按统一规则修正过度切分，使同一种表头含义对应稳定的结构。主要想法是：行、列、表头和跨格单元都是有明确边界的矩形，用目标检测就能预测它们，再由几何关系组装成逻辑表。

模型分为两个沿用DETR的检测器（\ref{sec:vision-recognition-detr}节）：一个在页面上定位整张表格，另一个在裁剪的表格图像中预测六类对象，即整表、行、列、列表头、投影行表头和跨格单元。\term{投影行表头}（projected row header）是横跨所有列、用于划分后续行组的标题行。两个检测器都以框和类别为监督，训练沿用DETR的二分匹配、类别交叉熵、$L_1$和广义IoU损失。数据集用于结构识别的部分含947642张表格，其中52.7\%至少含一个跨格单元。

两个检测器的结构相同，只有对象查询数不同，见表\ref{tab:vision-understanding-pubtables-config}。查询数取得略多于训练样本中对象数的最大值：表格检测取15个，结构识别取125个，因为一张表的行、列、表头和跨格对象远多于一页上的表格数。图像经ResNet-18主干得到步长32的特征图，用$1\times1$卷积降到256维并展平，加上位置编码后送入6层编码器；125个对象查询经6层解码器读取编码器输出，每个查询再由线性分类头输出6类对象或“无对象”，由三层感知机输出归一化的框。

\begin{table}[htbp]
  \centering
  \small
  \begin{tabularx}{\linewidth}{@{}p{14mm}XX@{}}
    \toprule
    部分 & 配置 & 输入与输出\\
    \midrule
    主干 & ResNet-18，ImageNet预训练，前几层冻结 & $H\times W\times3$ \ensuremath{\to} $H/32\times W/32\times512$\\
    投影 & $1\times1$卷积到256维，加位置编码 & 展平为$(HW/1024)\times256$\\
    编码器 & 6层，宽度256，8头，前馈2048，层归一化前置 & 序列长度不变\\
    解码器 & 6层；表格检测15个查询，结构识别125个查询 & $15\times256$或$125\times256$\\
    预测头 & 类别（结构识别为6类加“无对象”）与框 & 每个查询一组类别概率和框\\
    训练 & 20轮；结构识别初始学习率$5\times10^{-5}$，“无对象”类权重0.4，不用辅助解码损失 & —\\
    \bottomrule
  \end{tabularx}
  \caption{PubTables-1M所用DETR的配置。主干、层数、查询数、轮数、学习率与“无对象”权重见论文附录；宽度、头数、前馈维数、层归一化位置与辅助损失设置取自官方代码的结构识别配置文件。}
  \label{tab:vision-understanding-pubtables-config}
\end{table}

图\ref{fig:vision-table-structure}用四行三列的教学表说明组装过程。前两行为表头：“项目”跨两行，“金额（元）”跨第2、3列，其下分别为1月、2月。后两行是甲的$(12,18)$与乙的$(8,10)$。行列对象相交先产生12个基本格；两个跨格对象指出哪些基本格应合并，合并后形成10个逻辑单元；列表头对象标出前两行的功能。跨格表头使每个金额都能回溯到月份和单位。

\begin{figure*}[tbp]
  \centering
  % 构造表：4行3列；两个跨格对象各合并2格，得到10个逻辑单元。
\begin{tikzpicture}[x=1cm,y=1cm,font=\figurefont\small,
  arr/.style={-{Stealth[length=2mm]},BookInk,line width=0.9pt}]
  \node[anchor=west,font=\figurefont\bfseries\small] at (0,4.0)
    {行列交叉与跨格对象};
  \node[anchor=west,font=\figurefont\bfseries\small] at (9,4.0)
    {恢复后的逻辑表};
  \foreach \dx in {0,9}{
    \begin{scope}[xshift=\dx cm]
      \fill[BookBlue!12] (0,1.6) rectangle (2,3.2);
      \fill[BookTeal!12] (2,2.4) rectangle (6,3.2);
      \draw[BookInk,line width=0.9pt] (0,0) rectangle (6,3.2);
      \draw[BookRule,line width=0.5pt] (0,0.8) -- (6,0.8);
      \draw[BookRule,line width=0.5pt] (0,1.6) -- (6,1.6);
      \draw[BookRule,line width=0.5pt] (2,0) -- (2,3.2);
      \node[fill=BookBlue!12,inner sep=2pt] at (1,2.4) {项目};
      \node[fill=BookTeal!12,inner sep=2pt] at (4,2.8) {金额（元）};
      \foreach \x/\y/\v in {
        3/2/1月,5/2/2月,1/1.2/甲,3/1.2/12,5/1.2/18,
        1/.4/乙,3/.4/8,5/.4/10}
        \node at (\x,\y) {\v};
    \end{scope}
  }
  \draw[BookMuted,dashed,line width=0.5pt] (4,0) -- (4,3.2);
  \draw[BookMuted,dashed,line width=0.5pt] (0,2.4) -- (6,2.4);
  \draw[BookBlue,line width=1.1pt] (0,1.6) rectangle (2,3.2);
  \draw[BookTeal,line width=1.1pt] (2,2.4) rectangle (6,3.2);
  \node[fill=BookBlue!12,inner sep=2pt] at (1,2.4) {项目};
  \node[fill=BookTeal!12,inner sep=2pt] at (4,2.8) {金额（元）};
  \draw[BookRule,line width=0.5pt] (13,0) -- (13,2.4);
  \draw[BookRule,line width=0.5pt] (11,2.4) -- (15,2.4);
  \draw[arr] (6.3,1.6) -- (8.6,1.6);
  \node[align=center] at (7.45,2.35) {按跨格对象\\合并};
  \node at (3,-0.6) {12个基本格};
  \node at (12,-0.6) {10个逻辑单元};
\end{tikzpicture}

  \caption{行列检测到逻辑单元的教学恢复。左侧虚线表示行列交叉形成的基本格；右侧根据两处跨格对象合并表头。表头与单元归属是结构成果，字符内容来自另行读取；图中的数值供后续问答复算。}
  \label{fig:vision-table-structure}
\end{figure*}
\FloatBarrier

推理还要解决结构冲突，例如两个跨格对象同时声称拥有同一基本格。原实现通过抑制冲突对象或调整范围，使行、列及跨格关系可以组成逻辑表；之后把OCR或PDF中带位置的文字分配给单元，输出行列索引、跨度、内容和来源框。不能只输出检测框就声称已经恢复了HTML表格。

表格结构正确与字符串正确应分别评价。若甲2月的“18”被读成“13”，结构可以完全正确而后续差值错误；若“金额”被当作只覆盖1月，所有字符正确也仍损失单位归属。原工作不覆盖跨页表格，连续两页的表头重复和行连接需要额外处理。

\subsection{表单字段的解析}
\label{sec:vision-understanding-fields}

票据的“编号”“日期”“总额”不是规则的二维表格，而是一组键与值。输入是页面图像及其OCR结果，输出是字段记录：每个字段的名称、原文、规范化值和来源框。难点在于键与值的空间关系随版式变化，同一页上还可能出现多个相似金额。

常用做法分两步。第一步复用LayoutLMv3（\ref{sec:vision-understanding-layoutlm}节）的页面表示，按文字段预测角色，例如“问题”（即键）、“答案”（即值）、“标题”和“其他”，这需要相应的字段监督并微调任务头。第二步连接键与值：可以在成对文字段的表示上另训一个关系头，也可以使用清楚写出的规则；后者不能无条件归为LayoutLMv3原有结构。

设教学票据的读取结果含“编号A017”“日期2026-04-03”“已付100.00”“总额128.50”。一种简单规则是在每个键的同行右侧寻找值，要求日期符合日期格式、金额具有数值格式，并避免同一值被互不相容的字段重复占用。这条规则可以在该版式下恢复记录，但遇到“值在键下方”的模板就需调整；学到的关系头能适应更多版式，却需要成对的连接标注。

字段结果应保留原文字、规范化值和区域来源。例如，原文“128.50”对应数值128.50及单位元；不能因另有“已付100.00”就将两者都归为总额。

\subsection{图表数据的解析}
\label{sec:vision-understanding-charts}
\label{sec:vision-understanding-chart-reading}

图表用位置、长度、颜色或角度表达数值。输入是图表图像，输出是数据表：每个数值连同它所属的系列和类别。读取图表需要反向还原这些编码，再把数值对应到标签和系列；只恢复一个高度而不知道它属于哪个月份，仍不能回答月份比较。

ChartQA使用的数据表恢复管线在ChartOCR基础上扩展\scite{vision-chartqa2022}。它先用关键点检测网络找出绘图区、条形端点、折线顶点及文字区域，用CRAFT（\ref{sec:vision-understanding-craft}节）检测并读取轴标签和图例；条形由左上和右下关键点恢复，折线由连接点恢复，饼图由扇区角度恢复。随后依据纵轴标签估计刻度，再用颜色和位置邻近关系把数值关联到横轴标签与图例。这一管线由检测、读取和几何规则组成，数值换算本身不需要学习。

在线性纵轴上，已知像素纵坐标$y_0,y_1$分别对应数值$v_0,v_1$，则位置$y$的读数为
\begin{equation}
  v(y)=v_0+\frac{y-y_0}{y_1-y_0}(v_1-v_0)\text{，}
  \qquad y_1\neq y_0\text{。}
  \label{eq:vision-chart-calibration}
\end{equation}
这里纵坐标向下增加，因此较大的数值通常对应较小的$y$。公式来自两个锚点间的线性插值，不能用于未变换的对数轴。

沿用表格中甲的金额，设图表0元刻度位于$y_0=340$，20元刻度位于$y_1=140$，两根柱顶分别为220和160像素。代入得到$0+\frac{220-340}{140-340}\times20=12$元和$0+\frac{160-340}{140-340}\times20=18$元。检测结果还需依据横轴位置对应1月、2月，再依据图例把它们归到甲系列。图\ref{fig:vision-chart-reading}突出这个“像素—刻度—记录”的联系。

\begin{figure}[htbp]
  \centering
  % 数值与table-structure-recovery相同；像素轴 y=340-10v。
\begin{tikzpicture}[x=1cm,y=1cm,font=\figurefont\small,
  arr/.style={-{Stealth[length=2mm]},BookInk,line width=0.8pt}]
  \node[anchor=west] at (0,4.8) {金额（元）};
  \draw[arr] (0.8,0) -- (0.8,4.4);
  \draw[arr] (0.8,0) -- (5.0,0);
  \foreach \y/\v in {0/0,2/10,4/20}{
    \draw[BookMuted,line width=0.5pt] (0.7,\y) -- (0.8,\y);
    \node[anchor=east] at (0.65,\y) {\v};
    \draw[BookRule,dashed,line width=0.5pt] (0.8,\y) -- (4.7,\y);
  }
  \draw[BookBlue,line width=0.9pt,fill=BookBlue!18]
    (1.65,0) rectangle (2.35,2.4);
  \draw[BookBlue,line width=0.9pt,fill=BookBlue!18]
    (3.35,0) rectangle (4.05,3.6);
  \node[anchor=south] at (2,2.5) {$y=220$};
  \node[anchor=south] at (3.7,3.7) {$y=160$};
  \node[anchor=north] at (2,-0.1) {1月};
  \node[anchor=north] at (3.7,-0.1) {2月};
  \node[anchor=north] at (2.9,-0.7) {甲系列};
  \node[anchor=west,font=\figurefont\bfseries\small] at (5.5,3.7)
    {恢复记录};
  \node[anchor=west] at (5.5,2.7) {甲，1月，12元};
  \node[anchor=west] at (5.5,1.7) {甲，2月，18元};
  \node[anchor=west,align=left] at (5.5,0.5)
    {像素纵坐标向下\\金额沿纵轴向上};
\end{tikzpicture}

  \caption{线性纵轴的教学读数。轴上标注金额，柱顶旁标注原图像素纵坐标；右侧记录同时保留系列、月份、数值与单位。两根柱的数值与图\ref{fig:vision-table-structure}中甲行一致。}
  \label{fig:vision-chart-reading}
\end{figure}
\FloatBarrier

若纵轴是对数刻度，应在线性关系成立的$\log v$上插值，再取指数恢复$v$；若柱从非零基线开始，只比较柱长又会改变读数。颜色近似、图例遮挡和折线交叉会造成系列归属错误，需要独立于数值误差检查。将“18”误读成“13”后，后续计算器仍能准确算出$13-12=1$，却无法得到真实差值6；算术正确不能修复此前丢失的视觉证据。

\subsection{图解关系的解析}
\label{sec:vision-understanding-diagrams}
\label{sec:vision-understanding-diagram-graph}

流程图中的箭头表示有向关系，位置相邻的两个框未必相连。输入是图解图像，输出是有向图：节点及其文字，边及其方向和标签。难点在于连线可能交叉、弯折或断开，节点之间的关系只能通过追踪线迹来确定，不能由相对位置推断。

一种教学组合方法先用区域检测取得节点框，在各框内识别文字；随后检测线段和箭头尖端，沿连续线迹追踪连接；再依据端点与节点边界的接触，将线迹归到源节点和目标节点。节点检测可用框监督，连接预测可用端点或关系监督；在规则线稿中，也可以明确定义阈值、线段连接和箭头形状规则而不训练模型。

设节点集合为$V=\{A,B,C,D\}$，目标输出为有向图$\mathcal G=(V,E)$。若两条连线$A\to D$和$B\to C$在中间交叉，交点处没有连接点标记，恢复结果仍应只有这两条边；按“距离最近”接线可能误造$A\to C$和$B\to D$。断线、遮盖或缺少箭头时，应保留方向未知或多个候选，不能仅由页面上下顺序决定方向。

输出记录可以写成“源节点A，边标签通过，目标节点D”，并带原始线迹和端点位置。结构检查包括是否遗漏边、方向是否反转、边标签是否属于当前连接。Pix2Struct在图解问答上的微调说明，生成式模型可以直接回答关于图解的问题，但这并不等于输出了可检查的关系图；需要显式图结构时，仍要有相应监督和验证。

\section{视觉描述与内容概括}
\label{sec:vision-understanding-description}

定位和读取保留具体证据，描述则从中选择一部分组成语言。输入是图像、区域或视频，输出是一段文字；难点在于选什么内容、选出的内容能否在观测中找到依据。语言解码沿用\ref{sec:vision-models-language-generation}节的条件生成，这里要补足的是内容如何被选择，以及选择后如何保持与来源一致。

本节按描述覆盖的对象展开。全局场景的描述概括整幅图，要在有限长度内挑出主要对象与关系；局部对象的描述只讲指定区域，要在放大细节的同时保留与周围的关系；事件与步骤的描述面对视频，要决定哪些重复和细节可以压缩，哪些次数与先后必须保留。覆盖范围越窄，描述对证据精度的要求越高。

\subsection{全局场景的描述}
\label{sec:vision-understanding-global-caption}

\term{图像描述}（image captioning）以整幅图为输入，生成一句或一段概括性的文字。街景短描述可能只需指出“几辆车停在路边”；详细描述则需要补充车辆、人物、招牌与相互位置。不同长度并不意味着观察事实不同，但会让更多属性和关系进入输出，也就增加了需要核验的内容。

生成每个词时，模型都要决定从图像的哪里取信息。两条路线对“取信息的单位”给出不同回答。基于检测区域注意力的描述先用目标检测器把图像拆成一组对象区域，再在生成每个词时对这些区域做注意力；视觉单位是对象，语言解码器是为描述任务专门训练的循环网络。基于视觉语言模型的描述用通用视觉编码器的图像块或查询表示作为视觉单位，接入预训练的大型语言模型，用指令控制描述的范围和长度。

\subsubsection{基于检测区域注意力的描述}
\label{sec:vision-understanding-bottomup}

早期带注意力的描述模型在卷积特征图的均匀网格上分配注意力。一个对象常被网格切成若干块，模型很难把注意力恰好放在整个对象上；而描述中的名词大多对应完整对象。Bottom-Up/Top-Down的主要想法是把注意力的候选从网格单元换成检测器提出的对象区域：“自底向上”指区域由图像检测产生，“自顶向下”指当前语言状态决定读取哪些区域\scite{vision-bottomup2018}。

自底向上的部分是一个Faster R-CNN检测器，骨干为在ImageNet上预训练的ResNet-101，再在Visual Genome上训练预测1600类对象和400类属性；属性预测迫使区域特征保留颜色、材质等细节。候选框的非极大值抑制取IoU阈值0.7，各类别检测结果的抑制取0.3；保留任一类别置信度超过0.2的区域，每幅图最多100个，区域数随图像复杂程度变化。原论文还报告，初步实验中直接取前36个区域效果几乎相同。每个区域的特征经平均池化为2048维向量$\bm v_i$。描述模型训练期间，这些图像特征固定不变。

经典检测器（例如\ref{sec:vision-features-dpm}节的DPM）同样先确定候选窗口、再逐个判断类别；这里的差别在于区域特征是为对象和属性分类联合学习的深层特征，并直接作为后续语言模型的视觉单位，而不只是输出一个类别标签。从\ref{chap:vision-multimodal-alignment}章的角度看，检测器把注意力的候选限制在少数对象区域上，相当于一种“硬”注意力；后面的语言模型再在这些区域上做软注意力。

自顶向下的部分由两层LSTM组成，LSTM的门控结构可参考\BookChapterRef{chap:recurrent-networks}{卷五的循环神经网络章}。第一层注意力LSTM读取三样输入：前一时刻语言LSTM的输出、所有区域的平均特征，以及前一个词的嵌入，形成当前状态$\bm h_t^{\mathrm{att}}$。它据此计算
\begin{equation}
  \begin{aligned}
    e_{i,t}&=\bm w^\top\tanh(\bm W_{\mathrm v}\bm v_i
                             +\bm W_{\mathrm h}\bm h_t^{\mathrm{att}}),\\
    \alpha_{i,t}&=\frac{\exp(e_{i,t})}{\sum_{k=1}^{K}\exp(e_{k,t})},
    \qquad
    \hat{\bm v}_t=\sum_{i=1}^{K}\alpha_{i,t}\bm v_i\text{。}
  \end{aligned}
  \label{eq:vision-caption-attention}
\end{equation}
各区域特征与注意力状态先映射到相同维数，$\bm w$再把每个组合压成标量。归一化权重形成当前视觉信息$\hat{\bm v}_t$，它与注意力状态一起输入第二层语言LSTM；后者的输出经线性层和softmax预测下一个词。这是\ref{sec:vision-alignment-attention}节软对齐的直接应用：查询是语言状态，键和值是区域特征；区别只在打分形式，这里用一个带tanh的小网络计算分数，而不是缩放点积。

原论文给出了各层的宽度，表\ref{tab:vision-understanding-updown-config}按数据流列出每层的输入与输出维数。两层LSTM的隐藏单元数$M$都是1000，注意力层的隐藏维数$H$为512，词嵌入维数$E$为1000；词表由COCO训练描述中出现至少五次的词构成，共10010个。

\begin{table}[htbp]
  \centering
  \small
  \begin{tabularx}{\linewidth}{@{}p{18mm}XXp{14mm}@{}}
    \toprule
    层 & 输入 & 计算 & 输出维数\\
    \midrule
    区域特征 & 图像 & Faster R-CNN（ResNet-101），平均池化 & $K\times2048$\\
    注意力LSTM & 上一时刻语言LSTM输出（1000）、区域平均特征（2048）、上一词嵌入（1000），拼接为4048维 & LSTM，$M=1000$ & 1000\\
    注意力 & 区域特征与注意力LSTM状态 & 式\eqref{eq:vision-caption-attention}，$\bm W_{\mathrm v}\in\mathbb R^{512\times2048}$，$\bm W_{\mathrm h}\in\mathbb R^{512\times1000}$，$\bm w\in\mathbb R^{512}$ & 2048\\
    语言LSTM & 加权视觉特征（2048）与注意力LSTM输出（1000），拼接为3048维 & LSTM，$M=1000$ & 1000\\
    输出层 & 语言LSTM输出 & 线性层与softmax & 10010\\
    \bottomrule
  \end{tabularx}
  \caption{Up-Down描述模型的逐层维数。宽度取自原论文的实现细节；拼接维数由各输入维数相加得到。}
  \label{tab:vision-understanding-updown-config}
\end{table}

对于“一个人走过停着的汽车”，预测人物相关词时可以更多读取人物区域，描述汽车时再读取车辆区域。若两区域的打分为$(\log3,0)$，权重就是$(3/4,1/4)$；这表示当前计算中两种特征的贡献，并不是“这个词有75\%的概率由该区域证明”。

训练分两个阶段。先以图像及参考描述最小化逐词交叉熵，更新两个LSTM、注意力和输出层；原设置的学习率从0.01线性降到0，共6万次迭代，批大小100。随后用\term{自评序列训练}（self-critical sequence training）直接优化描述评价指标CIDEr：对同一幅图，采样生成一句$\bm w^{\mathrm s}$，再用贪心解码生成一句$\hat{\bm w}$作为基线，两句的CIDEr分数分别为$r(\bm w^{\mathrm s})$与$r(\hat{\bm w})$，参数沿
\begin{equation}
  -\bigl(r(\bm w^{\mathrm s})-r(\hat{\bm w})\bigr)\nabla_{\bm\theta}\log p_{\bm\theta}(\bm w^{\mathrm s})
  \label{eq:vision-scst}
\end{equation}
的方向下降。采样句比贪心句得分高时，提高它的概率，反之降低。这样不必让评价指标可导，也能直接改善它；代价是模型可能学会迎合指标偏好的措辞。CIDEr的定义见式\eqref{eq:vision-understanding-cider}；它对每句采样都能给出一个分数，正好用作这里的回报。使用时依赖已生成的词前缀逐步选择区域并以束搜索解码，原设置的束宽为5，并禁止同一个词连续出现两次。检测器遗漏的对象难以通过区域路径被描述，固定词表和训练分布又限制了可以表达的内容。注意力图能展示计算读取了哪里，但不能独立证明每个生成词都正确。

\subsubsection{基于视觉语言模型的描述}
\label{sec:vision-understanding-vlm-caption}

BLIP-2、Qwen3-VL等模型把视觉表示接入已有语言模型，连接方式见\ref{sec:vision-models-blip2}节和\ref{sec:vision-models-qwen3vl}节。它们的共同结论是：视觉编码器输出的表示经连接模块投影为语言模型可读的输入，语言模型在这些输入和文字提示的条件下逐词生成，预训练使其具备通用的图像描述和指令跟随能力。

用于描述时，输入由图像和描述指令组成。可以直接使用已有检查点，按其支持的提示格式要求“一句话概括”或“详细描述”；也可以在目标描述数据上微调。BLIP-2的描述微调是一个例子：语言模型保持冻结，以“a photo of”作为生成的起始提示，更新Q-Former和图像编码器，目标是参考描述的逐词交叉熵。描述指令改变输出选择，并没有增加图像中的事实。

同一张照片要求“一句话概括”时，可以覆盖主体活动；要求“说明右侧两辆车的差别”时，必须读取局部颜色、形状和相对位置。生成后可以将内容拆成“存在两辆车”“右车为蓝色”“左车位于前方”等判断，逐项回查。若某一细节因分辨率不足不可见，应缩小描述范围，或补充局部观测。这类模型的典型失败是\ref{sec:vision-understanding-challenge-language-prior}节所说的视觉幻觉，即输出了图像中没有依据的对象或属性；语言模型的先验越强，越可能补出“街景中通常会有”的内容。

两条路线在同一街景描述任务上的差别见表\ref{tab:vision-understanding-caption-compare}。视觉语言模型把检测区域换成通用的图像块或查询表示，把专门训练的LSTM换成预训练语言模型，换来了开放词汇、可由指令控制的长度与范围，以及不依赖检测器类别表的能力；代价是计算量大得多，视觉单位不再对应可检查的对象框，生成内容与图像区域的对应也更难追溯。

\begin{table}[htbp]
  \centering
  \small
  \begin{tabularx}{\linewidth}{@{}p{14mm}XX@{}}
    \toprule
    比较项 & 基于检测区域注意力的描述 & 基于视觉语言模型的描述\\
    \midrule
    视觉单位 & 检测器给出的对象区域 & 图像块或查询表示\\
    语言解码 & 专门训练的两层LSTM & 预训练大型语言模型\\
    训练 & 描述数据上交叉熵，再优化CIDEr & 大规模预训练，可选描述微调\\
    输出控制 & 由训练数据决定的固定风格 & 指令控制长度与范围\\
    可追溯性 & 注意力权重指向对象区域 & 需另行回查区域\\
    主要失败 & 检测遗漏、词表受限 & 视觉幻觉、细节因压缩丢失\\
    \bottomrule
  \end{tabularx}
  \caption{两种全局描述路线的对照。两者输入同一幅图，输出一段描述；评价都需要同时检查覆盖与事实性。}
  \label{tab:vision-understanding-caption-compare}
\end{table}

流畅度、覆盖和事实性是不同维度。短句可以忠实却遗漏关键内容，长句可以覆盖广却加入不存在的对象。评价时先固定期望描述范围，再判断哪些必要信息被覆盖、哪些输出得不到图像支持。

\subsection{局部对象的描述}
\label{sec:vision-understanding-region-caption}
\label{sec:vision-understanding-crop-caption}

全局压缩常弱化小物体细节。局部描述的输入是一幅图和一个指定区域，输出是关于该区域对象的文字。它可以提高目标区域的有效分辨率，但对象与参照物的关系可能跨越裁剪边界，因此需要同时保留来源与必要上下文。

一种组合流程先按查询取得区域框（\ref{sec:vision-understanding-grounding}节），裁剪后记录“源图编号、原图边界、缩放比例”，再将局部图和整图按模型支持的多图接口组织为输入。若模型本来接受框或区域标记，也可以使用其原生接口。训练区域描述器需要对应的“区域—描述”记录，不能把整图描述训练无条件当作已经学过所有区域指代。

假设两辆外观相近的车分别位于左右两侧，局部裁剪显示右车有车顶行李架；整图则显示它位于另一辆车后方。描述“右侧带行李架的车停在另一辆车后面”同时使用了局部外观和全局位置。若只保留右车裁剪，“在后面”的参照就被删掉了。结果应明确当前描述属于哪幅源图的哪个对象，不能把局部图中的左、右重新当作全图关系。

核验包括区域是否选对、局部属性是否读对、引用的关系是否仍有上下文支持。增加裁剪数量会增加读取成本，也可能产生相互矛盾的局部判断，因此应记录裁剪预算，避免只比较最终文字长度。

\subsection{事件与步骤的描述}
\label{sec:vision-understanding-event-caption}
\label{sec:vision-understanding-hierarchical-caption}

视频描述的输入是一段视频，输出是按时间组织的事件或步骤文字。概括会舍弃重复和细节，但事件次数与先后可能正是后续问题需要的信息。压缩的对象和保留条件应由请求决定，不能默认所有重复事件都可合并。

可以先用\ref{chap:vision-video-recognition}章的取样与片段表示，为每段保留实际时间，再交给支持视频或交错帧输入的语言接口形成描述，例如\ref{sec:vision-models-qwen3vl}节的接口。长视频也可以先生成局部描述，再按问题组织总体概括；这是一种分层组合，摘要应仍可追溯到片段。

设教学视频在$[1,2)$秒和$[5,6)$秒各发生一次敲门，$[3,4)$秒有人讲话。问“发生了什么”可以概括为“人物两次敲门，中间说了一句话”；问“第二次敲门是否在讲话之后”则必须保留第二次事件和讲话结束位置。只存储“人物敲门并讲话”会丢掉次数与顺序，即使该短描述本身没有虚构对象。

片段合并应避免把原顺序按主题重新排列，也不能把同一事件在相邻窗口中的重复描述当成两次事件。旁白、字幕和可见动作分别标明来源：字幕出现“我开门了”，不证明镜头中已经呈现开门。保存局部到总体的对应，可以在后续问答时重新取得被概括省略的证据。

\section{视觉问答与证据推理}
\label{sec:vision-understanding-qa}

\term{视觉问答}（visual question answering, VQA）依据图像、视频或文档等视觉内容回答问题。输入是视觉内容和一个问题，输出是答案，最好同时附带支持答案的证据。它可能只需读出一个日期，也可能需要连接多个时刻的事件。难点在于答案正确并不说明证据正确：模型可能凭问题中的语言规律猜对，也可能从错误的单元读出恰好相同的数。

前几节的检索、定位和读取可以成为问答的组成步骤。本节按回答一个问题的先后环节展开：先确定问题要求什么形式的答案；再按缺少的事实获取证据；需要时把多条证据组合起来；然后以选择、抽取、单元操作或生成的方式形成答案；最后核验答案与证据是否一致。

\subsection{答案要求的确定}
\label{sec:vision-understanding-question}

同一句“多少”可能要求不同的操作。存在问题要求判断是否有目标，属性问题要求给出颜色或状态，计数问题要求定义可数对象和范围，空间问题要求确定参照系，文档问题还需明确字段和单位。对“总额是多少”，应区分读取票据上已经写出的总额，还是依据若干金额自行求和；两种请求使用的操作不同，所需证据也不同。

问答模型容易利用问题中的语言规律猜测答案，例如“天空是什么颜色”多半答“蓝色”。VQA v2为同一问题收集答案不同的相似图像，使仅凭问题作答的捷径受到限制\scite{vision-vqav22017}。对同一问题更换图像后，答案是否随相关证据改变，是有用的诊断，但需要确认新图像确实改变了答案条件。

GQA利用场景图及问题程序构造组合问题，并检查一致性、视觉依据和答案合理性\scite{vision-gqa2019}。它能够系统改变属性、关系和组合深度，但这种构造也规定了题目分布，不能代表所有真实用户问题。面向连续对话时，还要保留图像编号、当前指代和已有回答的证据范围，不能只把此前答案作为不容质疑的事实。

\subsection{证据的获取}
\label{sec:vision-understanding-evidence}

确定答案要求之后，下一步是取得回答所缺的事实。获取证据应由当前缺少什么决定：读数问题需要足够分辨率，事件先后问题需要相应时间片段，多页文档问题则可能先要确定页面。每次增加输入，都应能够说明它补上了哪项信息。

证据所在的位置决定了获取方式，本小节按两种输入条件展开。页面与图像中的证据在空间上分布，获取靠页面检索、区域定位和局部放大，主要预算是分辨率和页数。长视频与多源输入中的证据在时间上分布，还可能分属画面、声音和字幕，获取靠时间检索、记忆压缩和跨来源对齐，主要预算是帧数和上下文长度。

\subsubsection{页面与图像中的证据获取}
\label{sec:vision-understanding-evidence-composition}

对“表格中甲的2月金额是多少”，若已给清晰单页，可以直接定位目标单元并读取；若给一份长文档，需要先找到包含该表的页面；若直接提供真实数据表，页面搜索和像素读取已由外部完成。三种条件的最终答案都可能是18元，但视觉工作量不同。

一种无需新增训练的组合规则是：先按标题或内容检索候选页（\ref{sec:vision-understanding-retrieval}节），再读取表格及其行列表头（\ref{sec:vision-understanding-tables}节），按“甲、2月”定位单元，最后返回数值、单位和来源。每条中间记录保留源文件或图像编号、页码、框或时间区间、原文字和处理方式。若目标页未被召回，应补充检索；若小数点看不清，应取得更清晰的裁剪。重复要求模型“仔细思考”并没有增加这两类缺失证据。

图\ref{fig:vision-qa-evidence}将同一教学表中的差值问题拆成可核验操作。这里的选择与计算是外显组合规则，各模型只承担已经说明的定位、读取或生成职责，不假设它们内部一定沿相同步骤推理。

\begin{figure}[htbp]
  \centering
  % 教学组合规则；不表示某个VLM内部必然具有该操作轨迹。
\begin{tikzpicture}[x=1cm,y=1cm,font=\figurefont\small,
  fact/.style={draw=BookBlue,fill=BookBlue!10,line width=0.9pt,
    minimum width=44mm,minimum height=13mm,align=center},
  operation/.style={draw=BookTeal,fill=BookTeal!10,line width=0.9pt,
    minimum width=50mm,minimum height=11mm,align=center},
  arr/.style={-{Stealth[length=2mm]},BookInk,line width=0.9pt,
    rounded corners=1mm}]
  \node[anchor=west] at (0,5.7) {问题：甲2月比1月多多少？};
  \node[fact] (jan) at (2.2,4.4)
    {甲，1月：12元\\页1，行3，列2};
  \node[fact] (feb) at (7.8,4.4)
    {甲，2月：18元\\页1，行3，列3};
  \node[operation] (subtract) at (5,2.25) {$18\text{元}-12\text{元}=6\text{元}$};
  \coordinate (join) at (5,3.2);
  \draw[BookInk,line width=0.9pt,rounded corners=1mm]
    (jan.south) -- (2.2,3.2) -- (join);
  \draw[BookInk,line width=0.9pt,rounded corners=1mm]
    (feb.south) -- (7.8,3.2) -- (join);
  \draw[arr] (join) -- (subtract.north);
  \node[operation,minimum width=60mm,fill=white] (answer) at (5,0.3)
    {甲2月比1月多6元};
  \draw[arr] (subtract.south) -- (answer.north);
\end{tikzpicture}

  \caption{带来源的表格问答教学流程。两个输入事实都指向同一页、同一行及不同月份单元；差值计算保留操作数与单位。任一读数或行列归属出错，都应返回相应证据检查。}
  \label{fig:vision-qa-evidence}
\end{figure}
\FloatBarrier

\subsubsection{长视频与多源输入中的证据获取}
\label{sec:vision-understanding-video-evidence}

视频问题需要的片段可能分散在很长的历史中。可以利用\ref{chap:vision-video-recognition}章的取样、检索和MovieChat式压缩记忆获得候选，再重新读取与问题相关的时刻。压缩历史保存了便于搜索的线索，并不保证保留每次手部动作或每个短暂出现的物体。

声音、字幕和图像应各自保留时间来源。画外有人说“放好了”，画面中的手却仍在移动物体，这可能来自说话者不同、语言含义不同或音画错位。支持音画输入的模型可采用\ref{sec:vision-models-qwen-omni}节的接口，具体区间对应可按\ref{sec:vision-understanding-av-time}节检查，不能只把不同来源文本拼在一起就认定它们同时描述同一事件。

LongVideoBench把视频与字幕共同作为较长上下文，要求利用所指片段的细节回答问题\scite{vision-longvideobench2024}。解释其结果时应分别记录帧数、视频时长、字幕及候选选项，避免把由字幕直接给出的答案全部归于视觉记忆。多选题的正确答案也没有同时验证空间框和精确时间区间。

两种条件下的获取都可以写成“缺什么、补什么”的循环，差别在于补充的代价。页面与图像中，补一次裁剪或一页的代价较小，主要风险是放大后丢掉上下文；长视频中，重新读取一段需要再解码视频帧，压缩记忆又可能已经丢掉所需细节，主要风险是证据在获取前就已不可恢复，以及多个来源在时间上错位。

\subsection{证据的多步组合}
\label{sec:vision-understanding-evidence-operations}

有些问题不能由单条证据回答，例如跨图比较数量，或判断两个事件的先后。输入是多条已核对的事实，输出是由它们推出的结论；难点在于每一步操作都要有明确的输入和依据，否则中间错误会被后续步骤掩盖。一种做法是把已核对的事实交给明确的操作规则执行，而不是要求模型在一次生成中隐式完成全部推理。

例如，两张图分别可见3辆和5辆符合条件的车，先分别保存对象编号和区域，再计算差值$5-3=2$。若问题问的是“新增了几辆”，只比较数量不足以回答，还需要跨图身份对应：增加2辆同时离开1辆，会使数量只增加1辆。

对前述敲门视频，可以将事件列表组织为“敲门1：$[1,2)$；讲话：$[3,4)$；敲门2：$[5,6)$”。由第二次敲门起点5大于讲话终点4，可回答“在讲话结束之后”。若两段存在重叠，则需要改用起点、终点或主要持续时段等更明确的关系，而不能强行给出唯一先后。

这种组合可以复用模型和计算工具，不需要虚构一个新的端到端训练目标。每步保留选取对象、读取事实、操作与结果，解释才有可核验的依据。NExT-QA通过时间关系、动作解释与场景问题考察更复杂的视频问答\scite{vision-nextqa2021}。其中“为什么”的标注属于对视频动作的语义解释，不等同于通过干预识别因果效应；仅凭“讲话之后敲门”无法证明讲话导致敲门。

\subsection{答案的形成}
\label{sec:vision-understanding-answering}

证据齐备后，还要把它们转成答案。输入是问题和证据的表示，输出是一个答案；难点在于答案的形式差别很大：可能是常见词“蓝色”，可能是只在当前图像中出现的地名，可能是文档原文中的一段日期，也可能是需要计算的差值。

不同的答案形式对应不同的输出计算，四种代表模型各采用一种。LoRRA从“固定答案词表加当前图像的OCR字符串”中选择一项，解决罕见词不在词表中的问题。DocVQA基线从文档的OCR文字序列中抽取一段连续跨度，答案保留原文拼写。VisionTaPas在数据表中选择单元并执行聚合操作，答案由所选单元算出。VL-T5逐词生成答案，输出不受候选集合限制。从前到后，输出空间越来越开放，而答案与证据之间的对应越来越隐蔽。

\subsubsection{LoRRA}
\label{sec:vision-understanding-lorra}

招牌上的罕见地名可能不在训练答案词表中。此前的VQA模型把问答当作在固定答案集合上的分类，词表外的答案永远无法输出，即使文字清清楚楚写在图里。TextVQA论文为此建立了需要读取图中文字的问答数据集，在OpenImages的28408幅图像上收集了45336个问题；并提出LoRRA，其主要想法是把当前图像的OCR字符串作为动态候选加入答案空间，使模型可以直接指向刚刚读到的内容\scite{vision-textvqa2019}。

模型有三个部分，结构见表\ref{tab:vision-understanding-lorra-config}。问题截取至多14个词，由GloVe词向量（词表取VQA 2.0问题中约7.7万个高频词）和带自注意力的LSTM编码为向量$\bm q$。视觉部分沿用VQA模型Pythia：网格特征取自预训练ResNet-152的res-5c块，经平均池化为2048维；区域特征取自在Visual Genome对象与属性上训练的Faster R-CNN的fc6层，训练时微调fc7层。两种特征在问题条件下分别做注意力加权平均，再拼接。读取部分先用外部OCR系统Rosetta取得至多50个OCR词，以FastText嵌入表示，使词表外的词也有向量；再以问题为条件对这些OCR词做注意力。外部OCR不与回答模块共同训练。论文没有给出注意力层和融合层的隐藏维数。

\begin{table}[htbp]
  \centering
  \small
  \begin{tabularx}{\linewidth}{@{}p{14mm}XX@{}}
    \toprule
    部分 & 输入与编码 & 输出\\
    \midrule
    问题 & 至多14个词；GloVe词向量，带自注意力的LSTM & 问题向量$\bm q$\\
    视觉 & ResNet-152 res-5c网格特征（2048维）与Faster R-CNN fc6区域特征；以$\bm q$为条件的注意力 & 两种加权特征拼接，与$\bm q$逐元素相乘\\
    读取 & Rosetta识别的至多50个OCR词；FastText嵌入；以$\bm q$为条件的注意力 & 加权特征与50个注意力权重拼接，与$\bm q$逐元素相乘\\
    答案 & 两路结果拼接；两层前馈网络 & $3996+50$个logit\\
    \bottomrule
  \end{tabularx}
  \caption{LoRRA的结构，依据TextVQA论文。固定答案词表含训练集中至少出现两次的3996个答案；隐藏维数论文未给出。}
  \label{tab:vision-understanding-lorra-config}
\end{table}

视觉与读取两路各自与问题表示融合，融合方式是逐元素相乘。按\ref{chap:vision-multimodal-alignment}章的分类，这是中间融合（\ref{sec:vision-alignment-intermediate-fusion}节）中最简单的乘性交互：问题向量的每一维像一个开关，放大或抑制对应维度的视觉特征，代价是只能表达同一维度之间的交互，不能像双线性融合那样让任意两维相互作用。两路融合结果拼接后送入答案模块。设固定答案数为$n$，当前OCR候选数为$m$，答案模块输出$n+m$个logit：前$n$个位置对应固定答案，后$m$个对应当前图像的OCR字符串，选中第$n+j$个位置就复制第$j$个字符串。读取部分的注意力只给出加权平均后的OCR特征，会丢失“哪个位置的词得了高分”这一信息；原实现因此把按原候选顺序排列的注意力权重也拼进特征，使答案模块知道应复制第几个候选。

训练使用逐候选的二元交叉熵，而不是要求唯一位置的softmax交叉熵。原因是同一答案可能同时出现在固定词表和OCR候选中，两个位置都应获得正监督。若固定词表没有“青山路”，OCR已正确读出它，模型仍可通过动态位置回答；若OCR把它切成“青山”和“路”两个候选，原LoRRA一次只能复制一个候选，不能输出多词答案；原论文统计，TextVQA中有8.9\%的问题只能由多个OCR词组合回答。OCR读错时，复制机制会原样复制错误。训练使用AdaMax优化器，批大小128，共2.4万次迭代，评价用\ref{sec:vision-understanding-challenges}节定义的VQA准确率。

\subsubsection{DocVQA基线}
\label{sec:vision-understanding-docvqa}

文档问答的答案往往已经出现在原文中，例如日期、编号或机构名。从固定词表中选择不适合这类答案，复制单个OCR词又不够表达多词答案。DocVQA数据集在12767页文档图像上收集了5万个问题，主要指标是ANLS（式\eqref{eq:vision-understanding-anls}）\scite{vision-docvqa2021}。论文同时测试了LoRRA和另一个文字问答模型M4C：原设置的LoRRA测试ANLS约为0.11，去掉固定答案词表后大幅下降，说明它主要依靠固定词表作答；M4C在动态OCR候选增加到500个时约为0.39，去掉检测器给出的对象特征后反而略有提高，说明为自然图像设计的对象区域在文档上用处不大。二者都不如下面的阅读理解基线。论文中表现最好的文字基线，主要想法是借用阅读理解中的\term{抽取式问答}（extractive question answering）：答案是上下文中的一段连续文字，模型只需预测这段文字的起点和终点，答案因此保留原有拼写。

基线把商业OCR软件识别出的词按从左上到右下的顺序、以空格连接成一个字符串作为上下文，与问题拼接后送入BERT编码。论文比较了三个预训练检查点：bert-base-uncased，以全词掩码预训练的bert-large，以及在其基础上又在SQuAD 1.1上微调过的bert-large-squad。BERT-base为12层、宽度768、12头，约1.1亿参数；BERT-large为24层、宽度1024、16头，约3.4亿参数\scite{devlin2019}。每层都是后置层归一化的标准Transformer编码器块。DocVQA论文没有给出最大序列长度和长文档的滑窗步长。设上下文词元表示为$\bm h_j$，两个线性头分别给出开始分数和结束分数，在全部上下文位置上做softmax，得到开始概率$p_{\mathrm s}(j)$和结束概率$p_{\mathrm e}(j)$。它利用的是OCR文字，不能因此被描述为BERT直接从像素理解布局。

对训练跨度$(s^*,e^*)$，目标为
\begin{equation}
  L_{\mathrm{span}}=-\log p_{\mathrm s}(s^*)-\log p_{\mathrm e}(e^*)\text{。}
  \label{eq:vision-docvqa-span}
\end{equation}
表现最好的配置采用已在SQuAD阅读理解数据上微调过的BERT-large，再在DocVQA上继续微调6轮，批大小8，学习率$2\times10^{-5}$；它在测试集上的ANLS为0.665，准确率为55.77\%。不经DocVQA微调、直接使用bert-large-squad时，测试ANLS只有0.475，说明阅读理解能力需要适应文档OCR文字的分布。数据只给答案字符串，原论文沿用TriviaQA的做法，用答案在序列中的首次匹配位置近似构造训练跨度；重复的日期或编号可能使这一弱标注指向错误位置。使用时在满足$s\leq e$和长度上限的组合中选择$p_{\mathrm s}(s)p_{\mathrm e}(e)$最大的跨度，再由词元到OCR字符串的映射恢复答案。

在教学票据中，序列包含“编号A017日期2026-04-03已付100.00总额128.50”，问题“票据日期是什么”应指向日期对应的跨度。若OCR已把03读成08，抽取模块即使选对跨度也会输出错误日期；若双栏顺序错乱，连续跨度又可能混入另一栏文字。要在结果中返回来源，可额外保留OCR词元到框的映射；LayoutLMv3等页面模型（\ref{sec:vision-understanding-layoutlm}节）以同样的跨度目标微调，只是把布局和图像也纳入了编码。

\subsubsection{VisionTaPas}
\label{sec:vision-understanding-visiontapas}

图表问题常需要计算，例如“甲的2月比1月多多少”。抽取式和选择式答案都无法表达一个不在页面上的差值。ChartQA数据集包含约九千六百个人工撰写的问题和约两万三千个由图表摘要生成的问题，其中人工问题大量涉及比较与算术\scite{vision-chartqa2022}。该论文提出的VisionTaPas的主要想法是沿用表格问答模型TaPas的做法：答案由“选哪些单元”和“对它们做什么聚合”两部分决定，模型只需预测这两部分，计算交给确定的规则；同时加入图表图像，使问题提到颜色或形状时也能利用视觉信息。

模型有三个编码部分，配置见表\ref{tab:vision-understanding-visiontapas-config}。图表图像切成图像块，经线性投影并加上一维可学习位置嵌入，送入12层的ViT-Base编码器，以ViT原论文发布的权重初始化\scite{dosovitskiy2020vit}；ChartQA论文没有给出输入分辨率和图像块大小。问题与按行展平的数据表以“[CLS] 问题 [SEP] 表格”的格式送入TaPas编码器。TaPas在BERT的基础上为每个词元加入所在行、列、数值排序以及上一轮答案等嵌入\scite{vision-tapas2020}，这里使用在WikiTQ表格问答数据上微调过的TaPas-Base权重；数据表可以是真实表，也可以由\ref{sec:vision-understanding-charts}节的管线从图表恢复。

两路输出进入四个跨模态块，每块包含视觉分支和文字表格分支。设第$\ell$块的视觉输入为$\bm H^{(\ell)}$，文字表格输入为$\bm Z^{(\ell)}$，视觉分支的计算顺序为
\[
 \begin{aligned}
  \bm H'&=\bm H^{(\ell)}+\operatorname{CA}(\bm H^{(\ell)},\bm Z^{(\ell)}),\\
  \bm H''&=\bm H'+\operatorname{MSA}(\bm H'),\\
  \bm H^{(\ell+1)}&=\bm H''+\operatorname{FFN}(\bm H'')\text{，}
 \end{aligned}
\]
其中$\operatorname{CA}(\bm A,\bm B)$表示以$\bm A$为查询、以$\bm B$为键和值的多头交叉注意力；每个子层都带层归一化。文字表格分支对称地以$\bm Z^{(\ell)}$为查询读取$\bm H^{(\ell)}$。两个分支的交叉注意力并行计算，读取的都是上一块的输出。这是\ref{sec:vision-alignment-intermediate-fusion}节中基于交叉注意力的中间融合：图像块可以根据表格中提到的系列名调整表示，表格单元也可以读取对应柱子的颜色和位置。末端在文字表格分支上接TaPas原有的两个头：单元选择头为每个单元给出被选概率，聚合头在若干操作中选择一个。

\begin{table}[htbp]
  \centering
  \small
  \begin{tabularx}{\linewidth}{@{}p{16mm}XX@{}}
    \toprule
    组件 & 配置 & 初始化\\
    \midrule
    视觉编码器 & ViT-Base，12层；分辨率与块大小论文未给出 & ViT原论文权重\\
    文字表格编码器 & TaPas-Base，12层，宽度768 & 在WikiTQ上微调的TaPas\\
    跨模态编码器 & 4块，每块两个分支：交叉注意力、自注意力、前馈层 & 论文未说明\\
    输出头 & 单元选择头；聚合头（无聚合、求和、计数、平均、差值、比值及固定答案类） & 前四种操作沿用TaPas\\
    训练 & ChartQA上30轮，学习率$10^{-5}$，批大小64 & —\\
    \bottomrule
  \end{tabularx}
  \caption{VisionTaPas的配置，依据ChartQA论文正文与附录。TaPas-Base的宽度沿用BERT-base；跨模态块的头数与宽度论文未给出。}
  \label{tab:vision-understanding-visiontapas-config}
\end{table}

原TaPas的聚合操作有无聚合、求和、计数和平均。ChartQA增加了差值和比值，并把“是”“否”等固定答案作为额外类别。设被选中的数值为$v_1,\ldots,v_k$，选择“求和”时返回$\sum_i v_i$，选择“计数”时返回$k$，选择“无聚合”时直接输出单元内容，选择“差值”时返回两个所选数值之差。对于“甲的2月金额比1月多多少”，正确选择是甲行的18和12，并选择差值，得到$18-12=6$元。单元选择指向具体来源，使人能够发现模型是否误用了乙行的10。

训练时，原TaPas可以只用最终答案作弱监督，让模型自己找到能算出答案的单元和操作。对新增的差值和比值，ChartQA改用启发式方法从答案和数据表直接推出候选单元：已知答案6，就在表内寻找差为6的数对。可能有多组数对满足同一差值，原论文人工检查一百例，发现约四分之一的这类监督有误。使用时选取概率超过阈值的单元，执行聚合头选出的操作。若请求“甲乙2月之和比1月之和多多少”，需要$(18+10)-(12+8)=8$元，这超出了单次聚合的表达，不能仅凭操作集中同时有“求和”和“差值”就认定支持嵌套计算。

\subsubsection{VL-T5}
\label{sec:vision-understanding-vlt5}

VisionTaPas的答案受限于操作集合，不能输出解释或任意格式的文字。VL-T5原本是一个把多种视觉语言任务统一为文字生成的模型：在T5编码器—解码器上，把图像区域特征和问题文字一起作为编码器输入，由解码器逐词生成答案\scite{vision-vlt5-2021}。ChartQA用它作为生成式基线，主要想法是不规定答案形式，让模型直接写出答案。

VL-T5以T5-Base为骨干\scite{raffel2020t5}，编码器和解码器各12层，宽度768，约2.2亿参数。原模型用在Visual Genome上训练的Faster R-CNN为每幅图取36个区域，每个区域的嵌入是四项之和：区域特征经线性层投影，框坐标经线性层投影，图像编号嵌入（多图输入时区分来自哪幅图），以及区域编号嵌入；区域编号嵌入与文字词表中新增的36个视觉哨兵词元共享，使模型可以用文字指称某个区域。文字词元与区域嵌入拼接成一个序列送入编码器，编码器块由自注意力和全连接层组成，解码器块另加交叉注意力。原模型在COCO与Visual Genome的约18万幅图像、918万个图文对上做多任务预训练。

文字输入由问题与序列化数据表组成。视觉输入沿用原VL-T5的接口，是一组对象区域特征及其框坐标；但自然图像检测器的对象类别不适合图表，ChartQA另以ResNet-101为骨干、以COCO权重初始化，训练了识别条形、折线、扇区、图例和坐标轴标签等15类图表元素的Mask R-CNN，提取这些元素的区域特征，数量不足时补零到固定的36个。编码器联合编码文字与区域特征，解码器根据多模态上下文及已有答案前缀预测后续词元。训练以目标答案的逐词交叉熵更新参数，在ChartQA上微调30轮，批大小96，学习率$10^{-4}$；使用时以束宽5的束搜索解码。评价采用\ref{sec:vision-understanding-challenges}节所说的放宽准确率。

对于前述差值问题，目标字符串可以是“6”；模型的输出空间不受某个聚合操作集合限制，但内部生成“6”的过程也不暴露18、12和减法。开放表达与准确计算是两种性质：生成接口可以输出完整解释，却仍可能在已经给对的数据上算错。判断视觉特征是否有贡献，还应固定数据表条件：若表头与数字足以独立回答，正确生成并不证明模型利用了图表图像；问题引用“蓝色柱”时，才可能需要颜色与系列的视觉对应。

四种答案形式的对照见表\ref{tab:vision-understanding-answer-compare}。从LoRRA到DocVQA基线，把“复制单个OCR词”换成“抽取连续跨度”，换来了多词答案，代价是依赖OCR序列化顺序。从抽取到VisionTaPas的单元选择与聚合，答案第一次可以是页面上没有的数，代价是只能表达预设的单步操作。VL-T5放开了输出形式，代价是答案与证据之间不再有显式对应。可追溯性与表达范围此消彼长，选择时应看任务更需要哪一项。

\begin{table}[htbp]
  \centering
  \small
  \begin{tabularx}{\linewidth}{@{}p{12mm}XXXX@{}}
    \toprule
    比较项 & LoRRA & DocVQA基线 & VisionTaPas & VL-T5\\
    \midrule
    答案形式 & 选择一项 & 连续跨度 & 单元加聚合 & 自由生成\\
    输出空间 & 固定词表与OCR词 & 上下文中的跨度 & 所选单元的运算结果 & 任意词元序列\\
    文字来源 & 外部OCR & 外部OCR & 数据表 & 数据表\\
    证据追溯 & 指向OCR候选 & 指向跨度位置 & 指向所选单元 & 无显式对应\\
    主要局限 & 只能复制一词 & 依赖序列化顺序 & 只有单步操作 & 计算过程不可见\\
    \bottomrule
  \end{tabularx}
  \caption{四种答案形式的对照。前两者原用于场景文字和文档问答，后两者用于同一ChartQA设置；比较时应固定文字来源与评价口径。}
  \label{tab:vision-understanding-answer-compare}
\end{table}

\subsection{答案的核验}
\label{sec:vision-understanding-answer-check}

答案形成后，还要确认它与证据和问题一致。答案可拆成对象、属性、字符串、数值、单位和关系等判断，分别与来源对照。若输出“总额是128.50美元”，原票据只写人民币符号，则数字正确而单位错误；若正确数值来自另一页的同名表格，来源也不成立。

核验可以改变相关观测后重新作答。例如，把甲2月金额从18改成13，差值应从6变成1；保持图像不变而改问乙行，答案应改为2元。前一种检查对输入变化的响应，后一种检查查询绑定，二者都需要重新确定正确答案，不能只观察模型是否“回答不同”。

缺少证据时回到相应获取步骤；来源冲突时保留冲突并限制结论。多轮对话中，“那一行”“他后来怎样”应绑定到已经明确的单元或对象，新增裁剪、页面或视频帧则登记为新的观测。上一轮回答本身不应替代尚未核对的原始事实。

\section{理解结果的核验}
\label{sec:vision-understanding-evaluation}

理解系统可能答对而找错证据，也可能位置正确但读错字符。验收要回答的问题是：最终结果之外，系统依据的中间结构、来源和输入是否都可靠。一个总分把这些情况混在一起，无法指出错误发生在哪里，也无法比较信息条件不同的系统。

本节按核验的维度展开，三个维度分别对应前文的不同层次。中间结构的核验检查检索候选、定位结果、字符串和表格等成果本身是否正确。来源证据的核验检查最终答案或描述是否真正由所报告的区域和时刻支持。模态贡献的核验检查答案依赖的是哪一种输入，以及比较的系统是否在相同的信息条件下工作。三者可以叠加使用，任一维度通过都不能代替其他维度。

\subsection{中间结构的核验}
\label{sec:vision-understanding-output-metrics}
\label{sec:vision-understanding-grounding-output}

匹配评价应固定候选单位，报告召回和排序；定位要检查区域、点或区间是否满足查询；读取和结构解析则应拆开字符、行列、阅读关系及连接错误。描述与问答还需核验事实覆盖和答案要求。表\ref{tab:vision-understanding-evaluation}给出沿本章例子的对应。

\begin{table}[htbp]
  \centering
  \small
  \begin{tabularx}{\linewidth}{@{}p{19mm}XX@{}}
    \toprule
    成果 & 直接检查 & 仍需单独检查\\
    \midrule
    检索候选 & 完整正例的召回和名次 & 关系是否由同一对象满足\\
    框或区间 & 与目标范围的吻合 & 指代、身份与空目标\\
    OCR字符串 & 字符编辑距离、整串正确 & 数字含义、单位与来源\\
    表格结构 & 行列、跨度及内容归属 & 单元数值和后续运算\\
    描述或答案 & 内容、数值与请求一致 & 支持它的区域和时刻\\
    \bottomrule
  \end{tabularx}
  \caption{多模态理解的分层核验。每行的直接检查只覆盖一种成果，不能自动证明最后一列的条件。}
  \label{tab:vision-understanding-evaluation}
\end{table}
\FloatBarrier

定位结果要能被复核。一个可复核的定位结果至少包含查询、源图或页面编号、时间信息、坐标约定及输出范围。框可以继续作为分割提示，点可以交给界面执行器，但新的输出要按新的标准检查。返回多个合理候选时，应说明差别来自参照不明、重复实例，还是目标部分遮挡；这种歧义反馈本身就是结果的一部分，不应被强行压成唯一答案。视频中“刚才拿杯子的人”还要求身份跨帧延续。可以复用\ref{sec:vision-video-tracking}节的对应和存在状态，区分进入、持续可见与离开；离开画面的对象不应被强行重新绑定到一个相似人物。2026年公开的TempoGround将跨帧对应、对象存在及二维和相机坐标系中的三维框联系起来，可作为这一问题的扩展案例\scite{vision-tempoground2026}。其二维定位、身份保持和三维几何仍是分别受约束的结果，论文设置不构成任意视频的长期一致性保证。

检索与定位中的关系绑定需要专门的诊断。Winoground固定两句文字的词汇集合，只改变词序，并要求与两幅图正确配对，可用于检查角色和关系绑定\scite{vision-winoground2022}。前述“人搬起箱子”和“箱子挡住人”也体现了相同诊断思路：不能只检测是否出现所有名词。基准论文中对当时模型的观察应限于所测设置，不外推为所有新模型的结论。

字符串的读取用字符错误率（式\eqref{eq:vision-cer}）衡量，并需事先规定空格、标点和大小写的处理。将“128.50”读成“123.50”只有一次替换，CER为$1/6$，但金额字段已经错误，因此业务字段还应报告整串或数值正确性。页面级解析则应分项报告文字、版面、表格和公式的误差，例如OmniDocBench的模块评价（\ref{sec:vision-understanding-page}节），并按页面类型保留结果，以发现报纸密集跨栏和论文公式等不同困难。

\subsection{来源证据的核验}
\label{sec:vision-understanding-grounded-evaluation}

答案正确不等于依据正确。来源证据的核验要求答案与它所报告的区域和时刻同时成立。若要求答案、时间和空间都正确，可以把单题成功写成三个条件的交集。令$A_i$表示答案正确，$T_i$表示返回时间区间满足阈值，$B_i$表示目标框满足阈值，则在明确标注协议下，联合成功率为
\begin{equation}
  \frac1n\sum_{i=1}^{n}\bm 1[A_i\land T_i\land B_i]\text{。}
  \label{eq:vision-grounded-success}
\end{equation}
这里$\bm1[\cdot]$为指示函数。这个教学指标的目的在于展示交集关系，不代称某个基准的全部计分规则。100题中若80题答案正确，其中50题时间也正确，这50题中又只有30题空间正确，联合成功率就是30\%，不是把三个环节的边际分数任意相乘。

2026年公开的VideoZeroBench通过不同层级分别收紧答案、时间和空间证据要求，具体体现了答案与依据可能分离的问题\scite{vision-videozerobench2026}。报告时还需说明模型是端到端寻找证据，还是已获得人工给定区间或框；额外给定证据改变了任务难度。该预印本的测试范围不能代表所有长视频应用。

描述和答案中不能出现无依据的内容。POPE以对象是否存在的询问检查视觉幻觉\scite{vision-pope2023}。应同时观察错误肯定和错误否定，避免总回答“没有”获得表面的保守。存在性正确仍不保证颜色、数量、动作、字符串或长段自由描述都正确，这些判断需要各自证据。空间判断还要保留观察视角：VSI-Bench考察视频中的空间信息保存与推断\scite{vision-vsibench2025}。视角改变后，“左侧”可能需要从当前相机坐标转换到场景关系；只记住一句“椅子在桌子左边”可能丢失当时参照。

% 完整作者边注不可跨页，为论述开头和来源共同预留阅读范围。
\Needspace{11\baselineskip}
解释长度也不是证据量。More Thinking, Less Seeing?研究了所测多模态推理模型中推理长度与视觉忠实性的关系，并提出相应诊断\scite{vision-morethinking2025}。由此需要检查的是延长输出后，新增判断是否仍有视觉依据；不能概括成“解释越长一定越好”或“一定越差”。若模型多写了五段推理，却始终使用读错的13，更多文字仍不能把正确差值恢复为6。

\Needspace{14\baselineskip}
Video-MME-v2用关联问题的成组评价补充单题准确率，使相互矛盾的回答更难被总分掩盖\scite{vision-videommev22026}。例如，模型既说“第二次敲门在讲话后”，又说“讲话发生在两次敲门之后”，在上述事件记录下至少一项错误。成组一致性是必要诊断，却不能独立证明内容正确：模型也可能一致地记错所有时间。其2026年版本应与所用字幕、音频和帧采样条件一起记录。

\subsection{模态贡献的核验}
\label{sec:vision-understanding-input-protocol}

多模态系统答对一个问题，未必用到了视觉。要判断视觉是否贡献了答案，可以移除、替换或错配图像，同时固定问题和其他输入。相同方法也适用于声音、字幕、OCR和数据表。移除模态可能产生训练分布之外的输入，性能变化不宜直接解释为纯粹因果贡献；多种受控替换与具体失败检查能提供更可靠的诊断。\ref{sec:vision-alignment-missing}节讨论过模型在推理时缺少某个模态怎样预测，\ref{sec:vision-alignment-imbalance}节讨论过联合训练中强势模态压制弱势模态的现象；这里的检查正是在评价端暴露这两种情形：若去掉图像后答案几乎不变，说明视觉在当前任务上几乎没有被利用，这与\ref{sec:vision-understanding-challenge-language-prior}节所说的语言捷径一致。

输入协议决定了系统实际承担的工作。若字幕已经明确写出答案，无图条件下答对并不意外。若真实数据表包含全部行列和数字，系统就不再承担从图像恢复表格的工作：ChartQA中同一模型在真实数据表和恢复数据表两种条件下的结果，衡量的是不同任务。报告需把原始像素、外部OCR、字幕、真实结构以及检索裁剪预算逐项固定，才能比较两个系统是否在相同信息条件下工作。

\section{本章小结}
\label{sec:vision-understanding-summary}

开篇的三个请求都要求把一句话落到特定观测上。本章按证据由粗到细的顺序给出了回答：检索先用双编码器从大库中召回候选，再以配对判定或逐条件核验重排；定位把文字短语、声音、音画事件和界面指令分别对应到区域、时间和点击位置；读取与解析恢复字符串，再按页面、表格、表单、图表和图解各自的结构组织关系；描述从中选择内容，问答则经过答案要求的确定、证据的获取与组合、答案的形成和核验。“右边拿杯子的人”由文字定位给出框，“总额是多少”由读取、字段解析和抽取给出带来源的数值，“第二次敲门在讲话之前还是之后”由时间定位与事件列表的比较给出。

第一节提出的五项挑战在各节分别得到处理。多粒度对应由召回、重排和四种定位分层解决；弱配对与噪声监督由间隔目标、困难负例和带置信度的伪标注缓解，但伪标注的自我强化仍是风险；高分辨率与长输入通过可变分辨率、裁剪搜索和按需获取证据处理，代价是更多的计算或调用；证据组合可以交给外显的操作或结构化的答案头，代价是表达范围受限；语言先验没有被某个模型消除，只能通过保留证据和分层核验暴露出来。贯穿各节的比较显示，同一个设计选择常同时带来收益和代价。双编码便于召回，却看不见关系绑定；联合计算能判断关系，却只能用于少量候选。无OCR的页面模型省去外部读取，却要求像素中保留小字；生成式答案表达自由，却不再暴露所依据的单元。可检查的结果应同时保留内容和来源：正确数字可能取自错误单元，正确描述可能遗漏关键时序，流畅解释也可能围绕错误观察展开。按中间结构、来源证据和模态贡献分层核验，才能确定问题发生在哪里。

仍未解决的问题集中在证据的完整性与组合推断上。关键证据可能在压缩或缩放时已经丢失，关系绑定错误与多步空间推断相互影响。Video-MSR在2026年将约束定位、链式指代、路线及反事实物理问题组织为多跳空间推理测试，可用于区分单步识别和组合推断\scite{vision-videomsr2026}。基准内训练带来的改善仍需放回具体任务评估，不能直接推论真实导航可靠性。Geo3R讨论了透视、物体朝向和视角变化引起的空间回答偏差，并在推理时引入几何信息及结构化三维关系\scite{vision-geo3r2026}。作为2026年公开案例，它说明可把空间判断重新落到几何对象上检查；若几何估计本身错误，增加三维表示也可能传播错误，而不是普遍消除视觉幻觉。

这些方向共同要求保留可核验的中间对象。读取越多并不总是越有用，关键在于缺少哪项事实、增加的观测是否补上它，以及不同来源冲突时怎样限定答案。下一章转向多模态生成：依据给定条件创作新内容时，需要保持的主体、空间和时间关系仍须明确，但未观测细节的生成与本章对既有证据的解释承担不同职责。

\paragraph{扩展阅读}
在大规模对比预训练普及之前，视觉语言理解多以检测区域特征加Transformer预训练的方式完成，可与本章基于检测区域注意力的描述模型和VL-T5对照阅读。ViLBERT用两路Transformer分别处理文字和区域，再以共同注意力层交换信息；LXMERT同样分设语言、对象关系和跨模态三个编码器，并加入掩码对象预测和图像问答等预训练任务；UNITER改为把文字和区域放入同一个Transformer，并用最优传输显式约束词与区域的对齐\scite{vision-vilbert2019,vision-lxmert2019,vision-uniter2020}OFA进一步把定位、描述、问答和图像生成都写成带指令的序列到序列任务，区域坐标也离散成词元输出，与VL-T5把任务统一为文字生成的思路一致，但覆盖的任务更广\scite{vision-ofa2022}

定位与读取方向，GLIP把目标检测改写成短语定位，用检测、定位和带伪标注的图文数据共同预训练，可与MDETR和Grounding DINO比较怎样扩展开放词表；Kosmos-2让多模态大语言模型在生成的文字中直接插入位置词元，把描述中的短语链接到图像区域\scite{vision-glip2022,vision-kosmos2-2023}Donut不调用外部OCR，由图像编码器和文本解码器直接从页面生成结构化输出，与本章先用TrOCR读取、再交给LayoutLMv3解析的流程正好相反，与Pix2Struct同属无OCR路线；Nougat沿用这种结构，把学术论文页面直接转换为包含公式和表格的标记文本\scite{vision-donut2022,vision-nougat2023}

问答与核验方向，OK-VQA收集了只靠图像内容无法回答、必须调用外部知识的问题，补充了VQA v2和GQA不考察的知识来源；EgoSchema用约三分钟的第一人称视频构造多项选择题，并以回答所需的最短视频长度衡量问题对长时证据的依赖，可与NExT-QA和LongVideoBench对照\scite{vision-okvqa2019,vision-egoschema2023}CHAIR在生成的图像描述中统计提到但图中并不存在的物体比例，是POPE之前常用的对象幻觉指标；POPE改用是非询问，两者分别从自由描述和受控提问两端暴露同一类错误\scite{vision-chair2018}
